\exercisesection{Regular rings}

\begin{enumerate}
\def\labelenumi{\arabic{enumi})}
\item
  Let A be a Noetherian local ring. For A to be regular, it is necessary and sufficient that we have \(\mathrm{di}_{\mathrm{A}}(\kappa_{\mathrm{A}})<+\infty\) .
\item
  Let A be a regular ring and p a prime ideal of A. Prove that the completion of A with respect to the p-adic topology is an integral domain.
\item
  Let A be a ring, and \(\mathbf{T} = (T_i)_{i\in I}\) a finite family of indeterminates. Prove that the following conditions are equivalent:
\end{enumerate}

\begin{enumerate}
\def\labelenumi{(\roman{enumi})}
\item
  the ring A is regular;
\item
  the ring A{[}T{]} is regular;
\item
  the ring A{[}{[}T{]}{]} is regular.
\end{enumerate}

\begin{enumerate}
\def\labelenumi{\arabic{enumi})}
\setcounter{enumi}{3}
\item
  Let A be a Noetherian ring of dimension \(\leqslant 2\) Show that for A to be regular, it is necessary and sufficient that the dual of every finitely generated A-module be projective (i.e., \(\mathfrak{p}\) a prime ideal of \(\Lambda\) ; if the second condition is satisfied, first prove that \(\mathrm{A}_{\mathfrak{p}}\) is regular when \(\operatorname{prof}(A_{\mathfrak{p}}) \leqslant 1\) and then that \(\mathrm{dp}_{\mathrm{A}_{\mathfrak{p}}}(\kappa(\mathfrak{p})) = 2\) when \(\operatorname{prof}(A_{\mathfrak{p}}) = 2\) ).
\item
  Let A be a Noetherian local ring of finite homological dimension; deduce from exerc. 16 of A, X, p.~208 another proof of the fact that A is regular (observe that one has dim \(_{\kappa_A}\) \(\mathfrak{m}_A/\mathfrak{m}_A^2 \leqslant \mathrm{dh}(A) \leqslant \mathrm{prof}(A)\) ).
\item
  Let A be a regular local ring, \((x_{1},\ldots,x_{n})\) and \((y_{1},\ldots,y_{n})\) maximal intersecting sequences of elements of \(m_{A}\) , generating ideals x and y respectively. We assume that we have \(y \subset x\) , so that there exist elements \(a_{ij}\) of A such that we have \(y_{i} = \sum_{j} a_{ij} x_{j}\) for \(1 \leqslant i \leqslant n\) . Prove that the class of \(\det(a_{ij})\) in A/y does not depend on the matrix \((a_{ij})\) chosen, and that its annihilator is the image of x; in particular, one has \(\det(a_{ij}) \neq 0\) (prove that the canonical homomorphism \(\operatorname{Ext}_{\mathrm{A}}^{n}(\mathrm{A}/x,\mathrm{A}) \to \operatorname{Ext}_{\mathrm{A}}^{n}(\Lambda/\mathfrak{y},\Lambda)\) is injective, then show using Koszul complexes that it identifies with the homomorphism A/x → A/y induced by the homothety of ratio \(\det(a_{ij})\) ).
\end{enumerate}

Let A be a regular local ring of dimension 3, and let f be a nonzero element of \(m_{A}\) Show that the following conditions are equivalent:

\begin{enumerate}
\def\labelenumi{(\roman{enumi})}
\item
  the ring A/(f) is not factorial;
\item
  there exists an integer \(n \geqslant 2\) and a matrix \(X \in \mathbf{M}_n(\mathrm{A})\) with elements in \(\mathfrak{m}_{\mathrm{A}}\) , such that \(f = \det(\mathrm{X})\) .
\end{enumerate}

(Under hypothesis (ii), prove that the ideal of A/(f) generated by the classes of the minors \(X^{11}, \ldots, X^{1n}\) of X, where \(X^{1p}\) is obtained by deleting the first row and the column of index p, has height 1 but is not principal. Under hypothesis (i), let p be a prime ideal of height 2 of A whose image in \(\Lambda/(f)\) is not principal; deduce from exerc. 11 of §3 that p is generated by the minors of order n-1 of a matrix \(Y \in \mathbf{M}_{n,n-1}(\mathbf{A})\) with elements in \(m_{\Lambda}\) . Write f as the determinant of a matrix \(\widetilde{Y}\) obtained by adding a column \(Y_{n}\) to Y; if an element of \(Y_{n}\) is invertible, reduce to the case where \(Y_{n}\) belongs to the canonical basis of \(\Lambda^{n}\) .

¶ 8) Let A be a regular local ring, \(\mathbf{x} = (x_1, \ldots, x_d)\) a coordinate system of A, \(\rho : A \to B\) an injective ring homomorphism making B a finitely generated A-module. Prove that the following conditions are equivalent:

\begin{enumerate}
\def\labelenumi{(\roman{enumi})}
\item
  for every integer \(p\) , the element \(\rho(x_1 \ldots x_d)^p\) does not belong to the ideal of B generated by \(\rho(x_1^{p+1}), \ldots, \rho(x_d^{p+1})\) ;
\item
  the A-module \(\rho(A)\) is a direct factor of B.
\end{enumerate}

These conditions are satisfied if A is a Q-algebra (exerc. 3 of § 2).

To prove (ii)⇒(i), adapt the proof of loc. cit. For the converse implication, reduce to the case where A is complete, and set, for \(p \geqslant 0\) \(A_{p} = A/(x_{1}^{p+1}, \ldots, x_{d}^{p+1})\) and \(B_{p} = B/(x_{1}^{p+1}, \ldots, x_{d}^{p+1})\) B. Observe that the class dc \((x_{1} \ldots x_{d})^{p}\) generates the socle of \(A_{p}\) and deduce from this that the homomorphism \(\rho_{p}: A_{p} \to B_{p}\) induced by \(\rho\) is injective. Prove that \(\rho_{p}\) admits a retraction \(A_{p}\) -linear. Conclude using E, III, p.~60, example II.

\begin{enumerate}
\def\labelenumi{\arabic{enumi})}
\setcounter{enumi}{8}
\tightlist
\item
  Let \(\rho: A \to B\) a local homomorphism of Noetherian local rings, \(N\) a finitely generated Macaulay B-module. Suppose that the ring \(A\) is regular and that we have \(\dim_{B}(N) = \dim(A) + \dim_{B/\mathfrak{m}_{A}B}(N/\mathfrak{m}_{A}N)\) . Prove that the A-module \(N\) is flat (argue by induction on \(\dim(A)\) ; if \(x \in \mathfrak{m}_{A} - \mathfrak{m}_{A}^{2}\) we will show that \(x\) is \(N\) -regular and that the homomorphism A/xA → B/xB and the (B/xB)-module N/xN satisfy the same hypotheses as ρ and N).
\end{enumerate}

¶ 10) Let A be a Noetherian ring, M a finitely generated A-module. For every integer \(n \geqslant 0\) , we denote by \(X_{n}\) the set of elements \(\mathfrak{p}\) of \(\operatorname{Spec}(A)\) such that \(\dim_{A_{\mathfrak{p}}}(M_{\mathfrak{p}}) - \operatorname{prof}_{A_{\mathfrak{p}}}(M_{\mathfrak{p}}) \geqslant n\) . Let \(k\) an integer \(\geqslant 0\) .

\begin{enumerate}
\def\labelenumi{\alph{enumi})}
\item
  For M to satisfy property \(S_{k}\) (§ 1, exerc. 8), it is necessary and sufficient that one have codim( \(X_{n}\) , Supp(M)) \(\geqslant n + k\) for every \(n \geqslant 0\) .
\item
  Assume that \(X_{n}\) is closed in \(\operatorname{Spec}(A)\) for every \(n \geqslant 0\) . Let p be an element of \(\operatorname{Supp}(M)\) . For the \(A_{p}\) -module \(M_{p}\) satisfies the property \(S_{k}\) (§ 1, exerc. 8), it is necessary and sufficient that one have \(\operatorname{codim}(X_{n}^{\prime}, \operatorname{Supp}(M)) \geqslant n + k\) for every \(n \geqslant 0\) and every component \(X_{n}^{\prime}\) of \(X_{n}\) containing p.
\item
  Assume that the ring A is presentable. Deduce from b) that the set of prime ideals p of A such that \(M_{p}\) satisfies the property \(S_{k}\) is open in Spec(A).
\end{enumerate}

\begin{enumerate}
\def\labelenumi{\arabic{enumi})}
\setcounter{enumi}{10}
\tightlist
\item
  Let A be the ring \(\mathbf{Z}[\mathbf{X}]\) , where \(\mathbf{X} = (X_i)_{i\in I}\) is a finite family of indeterminates. Let J be an ideal of A generated by monomials.
\end{enumerate}

\begin{enumerate}
\def\labelenumi{\alph{enumi})}
\item
  Let \(i \in I\) ; let \(\mathrm{P}_{1}, \ldots, \mathrm{P}_{r}\) ; \(\mathrm{Q}_{1}, \ldots, \mathrm{Q}_{s}\) of monomials generating J, where the \(\mathrm{P}_{j}\) are divisible by \(\mathrm{X}_{i}\) while c the \(\mathrm{Q}_{k}\) are not. Prove that the transporter J: \(\mathrm{X}_{i}\) (ideal of A consisting of the elements P such that \(\mathrm{X}_{i} \mathrm{P} \in \mathrm{J}\) ) is generated by the monomials \(\mathrm{X}_{i}^{-1} \mathrm{P}_{1}, \ldots, \mathrm{X}_{i}^{-1} \mathrm{P}_{r}\) ; \(\mathrm{Q}_{1}, \ldots, \mathrm{Q}_{s}\) .
\item
  Therefore deduce that if \(J: X_i\) is equal to \(J\) or \(A\) for every \(i \in I\) , \(J\) is the ideal generated by the variables \(X_i\) that belong to \(J\) .
\item
  Let M be an A-module, and let p be an integer. Assume that \(\operatorname{Tor}_{p}^{A}(A/J',M)\) (resp. \(\operatorname{Ext}_{A}^{p}(A/J',M)\) ) is zero for every ideal \(J'\) containing J and generated by some of the \(X_{i}\) . Prove that \(\operatorname{Tor}_{p}^{A}(A/J,M)\) (resp. \(\operatorname{Ext}_{A}^{p}(A/J,M)\) ) is zero (let \(J'\) an ideal containing J, generated by monomials, such that \(\operatorname{Tor}_{p}^{A}(A/J',M) \neq 0\) , and maximal for these properties; deduce from the exact sequence
\[
0 \rightarrow \mathrm{A} / (\mathrm{J} ^ {\prime}: \mathrm{X} _ {i}) \longrightarrow \mathrm{A} / \mathrm{J} ^ {\prime} \longrightarrow \mathrm{A} / (\mathrm{J} ^ {\prime} + \mathrm{AX} _ {i}) \rightarrow 0
\]

and hence b) that J' is generated by the elements \(X_{i}\) that it contains, which contradicts the hypothesis).

\item
  Hence deduce the inequality \(\mathrm{dp}_{\mathrm{A}}(\mathrm{A}/\mathrm{J}) \leqslant \mathrm{Card}(\mathrm{I})\) .

\item
  Let A be a ring, and \(\mathbf{x} = (x_{i})_{i \in I}\) a finite family of elements of A; suppose that for every subset I' of I the family \((x_{i})_{i \in I'}\) is completely secant for A. Let J be an ideal of A generated by elements \(x^{\alpha}\) , where \(\alpha\) ranges over a finite subset F of \(N^{1}\) . Prove the inequality dp \(_{A}\) (A/J) ≤ Card(I) (let J \(_{0}\) be the ideal of the ring A \(_{0}\) = Z{[}(X \(_{i}\) ) \(_{i \in I}\) {]} generated by the \(X^{\alpha}\) for \(\alpha \in F\) ; deduce from c) that Tor \(_{p}^{A_{0}}\) (A \(_{0}\) /J \(_{0}\) , A) is zero, and conclude by considering a projective resolution of the A \(_{0}\) -module A \(_{0}\) /J \(_{0}\) of length ≤ n).
\end{enumerate}

¶ 12) Let A be a regular local ring, M a finitely generated A-module, such that the A-module End \(_{A}\) (M) let it be free. We propose to prove that the following conditions are equivalent:

\begin{enumerate}
\def\labelenumi{(\roman{enumi})}
\item
  the A-module M is free;
\item
  the A-module M is reflexive;
\item
  \(\operatorname{Ext}_{\mathrm{A}}^{1}(\mathbf{M},\mathbf{M}) = 0\)
\end{enumerate}

\begin{enumerate}
\def\labelenumi{\alph{enumi})}
\item
  Prove the implication (iii) \(\Rightarrow\) (ii) (use exerc. 2 of VII, § 4).
\item
  Prove the converse implication (use Exerc. 4 if dim(A) ≤ 2, then reason by induction on dim(A) using Exerc. 18 of § 1).
\item
  Assume \(\dim(A) \leqslant 3\) Show the implication (ii) \(\Rightarrow\) (i) (use exerc. 4 if \(\dim(A) \leqslant 2\) ; if \(\dim(A) = 3\) , observe that (ii) implies \(\mathrm{dp}_{A}(M) \leqslant 1\) , then that (iii) and remark 2 of \(n^{\circ} 3\) imply that M is projective).
\item
  Assume M is reflexive. Let x be an element of \(m_{A} = m_{A}^{2}\) ; denote by B the ring A/xA, by N the B-module M/xM and by \(N^{*}\) its dual. Using b), prove that the B-module \(\operatorname{End}_{\mathbb{B}}(\mathbb{N})\) is free, then that \(\operatorname{End}_{\mathbb{B}}(\mathbb{N}^{*})\) is free (deduce from VII, § 4, no. 2 that \(\operatorname{End}_{\mathbb{B}}(\mathbb{N}^{*})\) is identified with the bidual of \(\operatorname{End}_{\mathbb{B}}(\mathbb{N})\) ).
\item
  Prove that (ii) implies (i) (argue by induction on the dimension of A, assumed ≥ 4; with the notation of d), the induction hypothesis implies that the B-module N* is free; deduce from this, using exerc. 18 of § 1 and prop. 7 of no. 4, that one has prof \(_{A}\) (Ext \(^{1}_{A}\) (M, B)) ≥ 1, then that Ext \(^{1}_{A}\) (M, A), which is of finite length by the induction hypothesis, is zero; conclude that the B-module M*/xM* is identified with N*, hence is free, then that the A-module M* is free and finally that M is free.)
\item
  Deduce from these results another proof of the fact that the ring A is factorial.
\end{enumerate}

\begin{enumerate}
\def\labelenumi{\arabic{enumi})}
\setcounter{enumi}{11}
\tightlist
\item
  Let A be a ring, and \((x_{1},\ldots,x_{n})\) a family of elements of A, generating an ideal I. We say that the family \((x_{1},\ldots,x_{n})\) is A-independent if the classes of \(x_{i}\) in I/I \(^{2}\) form a basis of this A/l-module.
\end{enumerate}

\begin{enumerate}
\def\labelenumi{\alph{enumi})}
\item
  Every completely intersecting family for A is A-independent.
\item
  Let \((x_{1},\ldots,x_{n})\) a A-independent family, with \(x_{1}=x_{1}^{\prime}x_{1}^{\prime\prime}\) . Prove that the families \((x_{1}^{\prime},\ldots,x_{n})\) and \((x_{1}^{\prime\prime},\ldots,x_{n})\) are A-independent, and that one has
\end{enumerate}

\[
\lg_ {\mathrm{A}} \left(\mathrm{A} / \left(x _ {1}, \dots , x _ {n}\right)\right) = \lg_ {\mathrm{A}} \left(\mathrm{A} / \left(x _ {1} ^ {\prime}, \dots , x _ {n}\right)\right) + \lg_ {\mathrm{A}} \left(\mathrm{A} / \left(x _ {1} ^ {\prime \prime}, \dots , x _ {n}\right)\right)
\]

(observe that the quotient \((x_{1}^{\prime},\ldots,x_{n})/(x_{1},\ldots,x_{n})\) is isomorphic to \(A/(x_{1}^{\prime\prime},\ldots,x_{n})\) ).

¶ 13) Let \(p\) be a prime number, \(q\) a power of \(p\) Let A be a reduced Noetherian local ring of characteristic \(p\) . One denotes by \(\mathrm{A}^q\) the subring of A consisting of the elements \(a^q\) for \(a \in \mathrm{A}\) . Prove that for A to be regular, it is necessary and sufficient that the \(\mathrm{A}^q\) -module A be flat.

(Reduce to the case where A is complete, hence a quotient of a formal power series algebra \(B = \kappa_{A}[[X_{1}, \ldots, X_{n}]]\) by an ideal b, which may be assumed to be contained in \(m_{B}^{2}\) . If A is regular, then b = 0. If A is flat over \(A^{q}\) let \(x_{i}\) be the image of \(X_{i}\) in A; prove that the family \((x_{1}^{q}, \ldots, x_{n}^{q})\) is A-independent (exercise 12), and deduce from this the equality \(\lg_{\mathrm{A}}(\mathrm{A}/(x_{1}^{q}, \ldots, x_{n}^{q})) = q^{n}\) . Show that this implies \(b \subset (X_{1}^{q}, \ldots, X_{n}^{q})\) ; prove similarly that one has \(b \subset (X_{1}^{q^{s}}, \ldots, X_{n}^{q^{s}})\) for every s, whence finally b = 0.)

\exercisesection{Complete intersections}

\begin{enumerate}
\def\labelenumi{\arabic{enumi})}
\item
  Let \(k\) a field, A the quotient ring of the polynomial ring \(k[X,Y,Z]\) by the ideal generated by \(Z^2,ZX,Z(Y-1)\) , \(x\) and \(y\) the classes of X and Y in \(\Lambda\) . Show that the sequence \((x,y)\) is completely secant for A but is not A-regular, although the sequence \((y,x)\) is A-regular.
\item
  Let A be a regular local ring, and I an ideal of A.
\end{enumerate}

\begin{enumerate}
\def\labelenumi{\alph{enumi})}
\tightlist
\item
  If ht(I) = 1, prove that the following conditions are equivalent:
\end{enumerate}

\begin{enumerate}
\def\labelenumi{(\roman{enumi})}
\item
  the ideal I is principal;
\item
  A/I is a Gorenstein ring;
\item
  A/I is a Macaulay ring.
\end{enumerate}

\begin{enumerate}
\def\labelenumi{\alph{enumi})}
\setcounter{enumi}{1}
\item
  Assume \(ht(I)=2\) for A/I to be a Gorenstein ring, it is necessary and sufficient that the ideal I be completely secant (if \((L,p)\) is a minimal projective resolution of I; observe that \(L_{1}\) is of rank 1 and consequently \(L_{0}\) of rank 2).
\item
  Let \(k\) a commutative field; we take \(A = k[[X,Y,Z]]\) and \(I = (XY,YZ,ZX, X^2 - Y^2, Y^2 - Z^2)\) . Prove that \(I\) is an ideal of height 3 that is not completely secant, and that \(A/I\) is a Gorenstein ring.
\end{enumerate}

\begin{enumerate}
\def\labelenumi{\arabic{enumi})}
\setcounter{enumi}{2}
\tightlist
\item
  Let R be a regular local ring, I an ideal of R, and A the ring R/I. Let \(\mathbf{x} = (x_{1}, \ldots, x_{d})\) a coordinate system of R, and \(\mathbf{y} = (y_{1}, \ldots, y_{m})\) a minimal generating system of I, so that \(m = [I/\mathfrak{m}_{R}\mathbf{l} : \kappa_{R}]\) .
\end{enumerate}

\begin{enumerate}
\def\labelenumi{\alph{enumi})}
\item
  When I is contained in \(\mathfrak{m}_{\mathrm{R}}^{2}\) define a canonical isomorphism of \(\mathrm{H}_1(\mathbf{x},\mathrm{A})\) on \(\mathrm{I} / \mathfrak{m}_{\mathrm{R}}\mathrm{I}\) (consider an exact sequence of Koszul complexes).
\item
  For the images of the \(x_{i}\) in A to form a minimal system of generators of \(m_{A}\) it is necessary and sufficient that 1 be contained in \(m_{R}^{2}\) .
\item
  Prove that the dimension of the \(\kappa_{\mathrm{A}}\) -vector space \(H_1(\mathbf{x}, A)\) is equal to \(m - d + [\mathfrak{m}_{\mathrm{A}} / \mathfrak{m}_{\mathrm{A}}^{2} : \kappa_{\mathrm{A}}]\) .
\end{enumerate}

4) Let A be a Noetherian local ring. Set \(\delta(A) = [\mathfrak{m}_{A}/\mathfrak{m}_{A}^{2}:\kappa_{A}] - \dim(A)\) (cf. § 4, no. 5).

\begin{enumerate}
\def\labelenumi{\alph{enumi})}
\item
  Show that the dimension of the \(\kappa_{\Lambda}\) -vector space \(H_{i}(\mathbf{z}, A)\) , where \(\mathbf{z}\) is a minimal generating system of \(\mathfrak{m}_{\mathrm{A}}\) is independent of the choice of \(\mathbf{z}\) ; it is denoted \(h_{i}(\mathrm{A})\) .
\item
  Prove the equality \(h_i(\mathrm{A}) = h_i(\widehat{\mathrm{A}})\) for every \(i\) .
\item
  Prove that one has \(h_1(A) \geqslant \delta(A)\) , and equality holds if and only if \(A\) is a complete intersection ring (Remark 2 of No. 2; reduce to the case where \(A\) is complete and apply Exerc. 4 by presenting \(A\) as a quotient \(R/I\) , where \(R\) is a regular local ring and \(I \subset m_R^2\) ).
\item
  Let \((x_{1},\ldots,x_{p})\) a completely secant sequence of elements of \(m_{A}\) generating an ideal J. Prove the equality \(h_{1}(A/J)-\delta(A/J)=h_{1}(A)-\delta(A)\) (reduce to the case p=1 and distinguish two cases, according as \(x_{1}\) belongs or does not belong to \(m_{A}^{2}\) ).
\end{enumerate}

\begin{enumerate}
\def\labelenumi{\arabic{enumi})}
\setcounter{enumi}{4}
\tightlist
\item
  Let R be a regular local ring, I an ideal of R, A the ring R/I. For A to be a complete intersection ring, it is necessary and sufficient that the ideal I be completely secant (use Exercises 4 c) and 3 c)).
\end{enumerate}

¶ 6) Let A be a Noetherian local ring, I an ideal of A.

\begin{enumerate}
\def\labelenumi{\alph{enumi})}
\tightlist
\item
  Prove that the following conditions are equivalent:
\end{enumerate}

\begin{enumerate}
\def\labelenumi{(\roman{enumi})}
\item
  the ideal I is completely secant;
\item
  the A/I-module I/I² is free and one has dp\_A(A/I) \textless{} +∞.
\end{enumerate}

Under hypothesis (ii), deduce from exerc. 2 of § 3 that there exists a cancellable element \(x\) in I such that \(x \notin \mathfrak{m}_{\mathrm{A}}\mathrm{I}\) ; reason by induction on the dimension of \(\mathrm{I} / \mathrm{I}^2\) .)

\begin{enumerate}
\def\labelenumi{\alph{enumi})}
\setcounter{enumi}{1}
\item
  Prove that the set of prime ideals p of A for which the ideal \(I_{p}\) is completely secant is open in Spec(A).
\item
  Let B be a presentable ring; prove that the set of prime ideals q of B such that \(B_{q}\) let a complete intersection ring be open in Spec(B).
\end{enumerate}

\begin{enumerate}
\def\labelenumi{\arabic{enumi})}
\setcounter{enumi}{6}
\tightlist
\item
  Let A be a ring, B an A-algebra; we assume that the A-module B is free of finite type. In the following exercises, we denote by B* the A-module Hom \(_{A}\) (B, A), and we endow it with its natural structure as a B-module.
\end{enumerate}

\begin{enumerate}
\def\labelenumi{\alph{enumi})}
\item
  Let \((e_i)_{i\in I}\) a basis of the A-module B, and \((e_i^*)_{i\in I}\) the dual basis of \(\mathbf{B}^*\) Prove that the element \(\mathrm{Tr}_{\mathrm{B / A}}\) of \(\mathbf{B}^*\) (A, III, p.~110) is equal to \(\sum_{j}e_{i}e_{i}^{*}\) .
\item
  We assume that the B-module \(B^{*}\) is free of rank 1; let \(\ell\) an element forming a basis of this module, and let \(\delta\) the element of B such that \(Tr_{B/A} = \delta\ell\) . The ideal \(\delta B\) of B does not depend on the choice of \(\ell\) ; it is called the different ideal of B over A. Prove that \(N_{B/A}(\delta)\) is a generator of the discriminant ideal of B over A (A, III, p. 115; compute the discriminant of a basis ( \(e_{i}\) ) of B on \(\Lambda\) using the matrix of multiplication by \(\delta\) in this basis).
\end{enumerate}

\begin{enumerate}
\def\labelenumi{\arabic{enumi})}
\setcounter{enumi}{7}
\tightlist
\item
  Let A be a ring, P a monic polynomial of degree \(n\) of A{[}X{]}, let B be the A-algebra A{[}X{]}/(P). We denote by \(x\) the class of X in B. Let \(\ell: \mathrm{B} \to \mathrm{A}\) be the A-linear form such that \(\ell(x^{n-1}) = 1\) , \(\ell(x^i) = 0\) for \(0 \leqslant i \leqslant n-2\) ; we denote \(\tilde{\ell}: \mathrm{B}[\mathrm{X}] \to \mathrm{A}[\mathrm{X}]\) the A-map{[}X{]}-linear map deduced from \(\ell\) by extension of scalars.
\end{enumerate}

\begin{enumerate}
\def\labelenumi{\alph{enumi})}
\item
  Let \(\sigma : B[X] \to A[X]\) the map \(A[X]\) -linear such that \(\sigma(x^i) = X^i\) for \(i = 0, \ldots, n - 1\) . Prove the formula \(P\tilde{\ell}(Q) = \sigma((X - x)Q)\) for every element \(Q\) of \(B[X]\) (check it for \(Q = 1, \ldots, x^{n-1}\) ).
\item
  Let \(\mathrm{P}^{*}(\mathrm{X}) = b_{n-1}\mathrm{X}^{n-1} + \ldots + b_{0}\) the polynomial B{[}X{]} such that one has \(\mathrm{P}(\mathrm{X}) = (\mathrm{X} - x)\mathrm{P}^{*}(\mathrm{X})\) in B{[}X{]}. Prove that \((b_{0}\ell, \ldots, b_{n-1}\ell)\) is the basis of \(B^{*}\) dual to the basis \((1, \ldots, x^{n-1})\) of B (apply a) to polynomials \(x^{i}\mathrm{P}^{*}\) ). Deduce that the B-module \(B^{*}\) is free of rank one, and that \(\ell\) forms a basis.
\item
  Prove the equality. \(\mathrm{Tr}_{\mathrm{B / A}} = \mathrm{P}'(x)\ell\) in \(\mathrm{B}^*\) (cf.~exerc. 7 a)).
\end{enumerate}

\begin{enumerate}
\def\labelenumi{\arabic{enumi})}
\setcounter{enumi}{8}
\tightlist
\item
  Let A be a ring, n an integer, S the algebra \(A[X_{1},\ldots,X_{n}]\) (resp. \(A[[X_{1},\ldots,X_{n}]]\) ), and a an ideal of S generated by a sequence \(\mathbf{P}=(\mathbf{P}_{1},\ldots,\mathbf{P}_{n})\) completely secant for S. We set \(B=S/a\) and we denote by \(\pi\) the canonical homomorphism from S onto S/a. Assume that the A-module B is free of finite type. We propose to prove that the B-module \(B^{*}\) is free of rank 1 and the different ideal of B over A is generated by \(\pi(\det(\frac{\partial P_{i}}{\partial X_{j}}))\) .
\end{enumerate}

\begin{enumerate}
\def\labelenumi{\alph{enumi})}
\item
  We set \(C = S \otimes_{A} B\) and \(\xi_{i} = X_{i} \otimes 1 - 1 \otimes \pi(X_{i})\) for \(1 \leqslant i \leqslant n\) . We denote by \(\varphi : C \to B\) the homomorphism of A-algebras such that \(\varphi(P \otimes b) = \pi(P)b\) . Prove that the kernel of \(\varphi\) is generated by the sequence \(\boldsymbol{\xi} = (\xi_{1}, \ldots, \xi_{n})\) , and that this sequence is completely secant for C.
\item
  Let \((c_{ij}) \in \mathbf{M}_{n}(\mathbf{C})\) be a matrix satisfying \(P_{i} \otimes 1 = \sum_{j} c_{ij} \xi_{j}\) . Prove the equality \(\pi(\frac{\partial P_{i}}{\partial X_{j}}) = \varphi(c_{ij})\) .
\item
  The matrix \((c_{ij})\) defines a morphism of complexes \(\mathbf{K}_{\bullet}(\mathbf{P},\mathrm{C})\to\mathbf{K}_{\bullet}(\boldsymbol{\xi},\mathrm{C})\) (A, X, p.~151), whence a morphism of complexes of S-modules \(u:\mathbf{K}_{\bullet}(\mathbf{P},\mathrm{S})\to\mathbf{K}_{\bullet}(\boldsymbol{\xi},\mathrm{C})\) The homomorphism \(u_{n}:S\to C\) maps \(1_{S}\) on the element \(d=\det(c_{ij})\) of C.
\item
  Let \(v: \mathrm{Hom}_{\mathrm{S}}(\mathrm{C}, \mathrm{S}) \otimes_{\mathrm{C}} \mathrm{B} \to \mathrm{B}\) be the B-linear homomorphism such that \(v(f \otimes 1) = \pi(f(d))\) ; prove that \(v\) is an isomorphism (identify \(v\) with \(\mathrm{H}^n(\mathrm{Homgr}_{\mathrm{S}}(u, 1_{\mathrm{S}}))\) , and observe that each of the complexes \(\mathbf{K}_{\bullet}(\mathbf{P}, \mathrm{S})\) and \(\mathbf{K}_{\bullet}(\boldsymbol{\xi}, \mathrm{C})\) defines a resolution of B by free S-modules of finite type).

\item
  We consider the composite map \(t: \mathrm{Hom}_{\mathrm{A}}(\mathrm{B}, \mathrm{A}) \to \mathrm{Hom}_{\mathrm{S}}(\mathrm{C}, \mathrm{S}) \to \mathrm{B}\) , where the first arrow is obtained by extension of scalars and the second is deduced from \(v\) . Prove that \(t\) is B-linear and bijective (identify \(\mathrm{Hom}_{\mathrm{S}}(\mathrm{C}, \mathrm{S})\) with \(\mathrm{Hom}_{\mathrm{A}}(\mathrm{B}, \mathrm{A}) \otimes_{\mathrm{B}} \mathrm{C}\) ).
\item
  Prove the formula \(t(\mathrm{Tr}_{\mathrm{B / A}}) = \varphi (d)\) (express \(d\) and \(\mathrm{Tr}_{\mathrm{B / A}}\) using a basis of B and the dual basis of \(\mathrm{B}^*\) cf. exercise 7, a)).
\item
  Deduce from b), e), and f) that the different ideal of B over A is generated by \(\pi(\det(\frac{\partial P_i}{\partial X_j}))\) .
\end{enumerate}

\begin{enumerate}
\def\labelenumi{\arabic{enumi})}
\setcounter{enumi}{9}
\tightlist
\item
  Let A be a discrete valuation ring, v its normalized valuation, and K its field of fractions. Consider a Noetherian A-algebra B equipped with a homomorphism of A-algebras \(\varepsilon:B\to A\) Let I denote the kernel of \(\varepsilon\) and J its annihilator in B.
\end{enumerate}

\begin{enumerate}
\def\labelenumi{\alph{enumi})}
\tightlist
\item
  Prove that the following conditions are equivalent:
\end{enumerate}

\begin{enumerate}
\def\labelenumi{(\roman{enumi})}
\item
  the \(\Lambda\) -module \(1 / I^2\) is of finite length;
\item
  the local ring of \(K \otimes_{A} B\) at the maximal ideal \(K \otimes_{A} I\) is isomorphic (as a K-algebra) to K ;
\item
  ; we have \(\mathrm{K}\otimes_{\mathrm{A}}\mathrm{B} = (\mathrm{K}\otimes_{\mathrm{A}}\mathrm{I})\oplus (\mathrm{K}\otimes_{\mathrm{A}}\mathrm{J})\)
\end{enumerate}

We now assume that these conditions are satisfied; we denote by \(\gamma(B, \varepsilon)\) , or simply \(\gamma(B)\) the length of the A-module \(I/I^{2}\) , and by \(\eta(B, \varepsilon)\) , or simply \(\eta(B)\) , that of the A-module \(A/\varepsilon(J)\) .

\begin{enumerate}
\def\labelenumi{\alph{enumi})}
\setcounter{enumi}{1}
\item
  Let b be an ideal of B contained in I; one considers the A-algebra \(B' = B/b\) and the homomorphism \(\varepsilon': B' \to A\) deduced from \(\varepsilon\) . Prove that \(B'\) also satisfies the conditions of a), and that one has \(\gamma(B', \varepsilon') \leqslant \gamma(B, \varepsilon)\) and \(\eta(B', \varepsilon') \leqslant \eta(B, \varepsilon)\) .
\item
  Assume that the A-module B is free of finite type. Prove that J is a free A-module of rank one, a direct factor in B, and that one has \(\mathrm{Tr}_{\mathrm{B / A}}(x) = \varepsilon (x)\) for every \(x\in J\)
\item
  Assume moreover that B is a Gorenstein ring; denote by \(\ell\) a basis of the B-module B* (§ 3, exerc. 5), and by δ a generator of the different ideal of B over A (exerc. 7, b)). Prove the equality η(B) = v(ε(δ)) (prove that \(\ell\)(J) is equal to A, and use c)).
\item
  One places oneself in the situation of \(b\) ), with \(\mathfrak{b} \neq 0\) ; assume that the A-algebras B and B' are finite and flat, and that B is a Gorenstein ring. Prove that one has \(\eta(B') < \eta(B)\) (observe that the image of J in B/ \(\mathfrak{m}_{A}\) B is the socle of B/ \(\mathfrak{m}_{A}\) B, hence is contained in \(\mathfrak{b}/\mathfrak{m}_{A}\mathfrak{b}\) ).
\end{enumerate}

¶ 11) One keeps the notation of the preceding exercise; assume that the rings A and B are local and complete.

\begin{enumerate}
\def\labelenumi{\alph{enumi})}
\item
  Show that B is isomorphic to the quotient of a formal power series algebra \(S = A[[X_{1}, \ldots, X_{n}]]\) by an ideal a contained in \((X_{1}, \ldots, X_{n})\) ; one may take \(n = \dim_{\kappa_{A}}(\mathfrak{n}/\mathfrak{n}^{2})\) , where n is the image of \(m_{B}\) in \(B/m_{A}B\) .
\item
  Prove that the following conditions are equivalent:
\end{enumerate}

\begin{enumerate}
\def\labelenumi{(\roman{enumi})}
\item
  B is a finite flat A-algebra and a complete intersection ring;
\item
  the ideal a is generated by a sequence \((\mathrm{P}_{1},\ldots,\mathrm{P}_{n})\) completely secant for \(\kappa_{A}\otimes_{A}S\) .
\end{enumerate}

(cf. exerc. 5.)

\begin{enumerate}
\def\labelenumi{\alph{enumi})}
\setcounter{enumi}{2}
\item
  If these conditions are satisfied, prove that one has \(\eta(B) = \gamma(B) = v(\varepsilon(\det(\frac{\partial P_i}{\partial X_j})))\) (use exercises 9 and 10, \(d\) )).
\item
  Show that one can find a sequence \((f_{1},\ldots,f_{n})\) of elements of a completely secant for \(\kappa_{A}\otimes_{A}S\) such that the A-algebra \(\mathrm{B}^{\prime}=\mathrm{S}/(f_{1},\ldots,f_{n})\) satisfies \(\gamma(\mathrm{B}^{\prime})=\gamma(\mathrm{B})\) .
\end{enumerate}

(Deduce from the exact sequence \(\mathfrak{a}/\mathfrak{a}^{2}\to\Omega_{\mathrm{A}}(\mathrm{S})\otimes_{\mathrm{S}}\mathrm{B}\to\Omega_{\mathrm{A}}(\mathrm{B})\to0\) (A, III, p.~137) an exact sequence of A-modules \(\mathfrak{a}/\mathfrak{a}^{2}\otimes_{\mathrm{S}}\mathrm{A}\to\Omega_{\mathrm{A}}(\mathrm{S})\otimes_{\mathrm{S}}\mathrm{A}\to\mathrm{I}/\mathrm{I}^{2}\to0\) ; take elements \(f_{1},\ldots,f_{n}\) of \(\mathfrak{a}\) whose images in \(\Omega_{\mathrm{A}}(\mathrm{S})\otimes_{\mathrm{S}}\Lambda\) form a basis over A of the image of \(\mathfrak{a}/\mathfrak{a}^{2}\otimes_{\mathrm{S}}\mathrm{A}\) .)

\begin{enumerate}
\def\labelenumi{\alph{enumi})}
\setcounter{enumi}{4}
\tightlist
\item
  Hence, using Exercise 10, e), we have \(\gamma(B) \geqslant \eta(B)\) and equality holds if and only if B is a complete intersection ring.
\end{enumerate}

\exercisesection{Extension of scalars in regular algebras}

\begin{enumerate}
\def\labelenumi{\arabic{enumi})}
\tightlist
\item
  Let \(k\) be a non-perfect field of characteristic \(p > 2\) , let \(a\) be an element of \(k - k^p\) , and let \(A = k[\mathrm{X},\mathrm{Y}] / (\mathrm{X}^2 -\mathrm{Y}^p +a)\) .
\end{enumerate}

\begin{enumerate}
\def\labelenumi{\alph{enumi})}
\item
  Prove that the ring A is regular (one may prove that A is integrally closed in its field of fractions \(K = k(Y)[X]/(X^{2} - Y^{p} + a)\) ).
\item
  Prove that k is algebraically closed in K.
\item
  Let \(k'\) the extension \(k(a^{1/p})\) of \(k\) prove that the ring \(A \otimes_k k'\) is not normal.
\end{enumerate}

\begin{enumerate}
\def\labelenumi{\arabic{enumi})}
\setcounter{enumi}{1}
\tightlist
\item
  Let \(k\) be a field, A an essentially finitely generated \(k\)-algebra, B a \(k\)-algebra, and \(p\) an integer.
\end{enumerate}

\begin{enumerate}
\def\labelenumi{\alph{enumi})}
\item
  If A and B satisfy the property \(S_{p}\) (§ 1, exerc. 8), then so does \(A \otimes_{k} B\) (argue as in the proof of prop. 4, using loc. cit., f)).
\item
  Let K be an extension of k; for \(\mathrm{A}_{(\mathbb{K})}\) check \(S_{p}\) it is necessary and sufficient that it be so for A.
\end{enumerate}

\begin{enumerate}
\def\labelenumi{\arabic{enumi})}
\setcounter{enumi}{2}
\tightlist
\item
  Let A be a Noetherian ring and \(\rho : A \to B\) a ring homomorphism. We say that the A-algebra B is absolutely regular if B is a Noetherian ring and a flat A-module and, for every \(\mathfrak{p} \in \operatorname{Spec}(A)\), the \(\kappa(\mathfrak{p}) \otimes_{A} B\) is absolutely regular.
\end{enumerate}

\begin{enumerate}
\def\labelenumi{\alph{enumi})}
\item
  Let T be a multiplicative subset of B and S a subset of A such that \(\rho(S) \subset T\) If B is a regular A-algebra, \(T^{1}B\) is an \(S^{-1}A\) regular algebra.
\item
  For a Noetherian A-algebra B to be absolutely regular, it is necessary and sufficient that for every maximal ideal n of B, the \(A_{\rho^{-1}(n)}\) -algebra \(B_{n}\) be absolutely regular.
\item
  If C is an absolutely regular B-algebra and B is an absolutely regular A-algebra, then C is an absolutely regular A-algebra.
\item
  If B is a Noetherian ring and C is a faithfully flat B-module and an absolutely regular A-algebra, then B is an absolutely regular A-algebra.
\item
  Let B be an absolutely regular A-algebra and A' an A-algebra essentially of finite type. The A'-algebra A' ⊗A B is absolutely regular.
\end{enumerate}

\begin{enumerate}
\def\labelenumi{\arabic{enumi})}
\setcounter{enumi}{3}
\item
  Let A be a Noetherian ring, B a faithfully flat absolutely regular A-algebra (exerc. 3). For A to be a regular (resp. normal, resp. Macaulay, resp. Gorenstein) ring, it is necessary and sufficient that the same be true of B.
\item
  One says that A is a Grothendieck ring if A is a Noetherian ring and, for every prime ideal p of A, the completion of the local ring \(\Lambda_{p}\) is an \(A_{p}\) -algebra absolutely regular (exerc. 3).
\end{enumerate}

\begin{enumerate}
\def\labelenumi{\alph{enumi})}
\item
  Prove that every ring of fractions and every quotient of a Grothendieck ring is a Grothendieck ring.
\item
  Let A be a Noetherian ring and B a faithfully flat absolutely regular A-algebra. If B is a Grothendieck ring, prove that the same is true of A (if p is a prime ideal of A and q a prime ideal of B above p, one will show that \(\widehat{B}_{q}\) is a \(\widehat{\Lambda_{p}}\) -module faithfully flat using III, § 5, n° 4, prop. 4; one will then apply exerc. 3).
\end{enumerate}

\begin{enumerate}
\def\labelenumi{\arabic{enumi})}
\setcounter{enumi}{5}
\item
  Let A be a Grothendieck ring, I an ideal of A, and \(\widehat{\mathsf{A}}\) the separated completion of A for the I-adic topology. Prove that \(\widehat{\mathsf{A}}\) is an absolutely regular A-algebra. (Let \(\mathfrak{n}\) be a maximal ideal of \(\widehat{\mathsf{A}}\) , and by \(\mathfrak{m}\) its inverse image in A; deduce from exerc. 3 that \(\widehat{\Lambda}_{\mathfrak{n}}\) is an \(\mathrm{A}_{\mathfrak{m}}\) -algebra absolutely regular.)
\item
  \begin{enumerate}
  \def\labelenumii{\alph{enumii})}
  \tightlist
  \item
    Show that an integral Dedekind ring of characteristic 0 is a Grothendieck ring.
  \end{enumerate}
\end{enumerate}

\begin{enumerate}
\def\labelenumi{\alph{enumi})}
\setcounter{enumi}{1}
\tightlist
\item
  Let \(k\) a field of characteristic \(p > 0\) such that \(k\) be of infinite dimension over \(k^p\) , \(r\) an integer \(\geqslant 1\) , and A the subring of \(k[[\mathrm{X}_1, \ldots, \mathrm{X}_r]]\) formed by formal series whose coefficients generate a vector space of finite dimension over \(k^p\) . Show that A is a regular local ring of dimension \(r\) whose completion is identified with \(k[[\mathrm{X}_1, \ldots, \mathrm{X}_r]]\) , but is not a Grothendieck ring (use exerc. 17 of III, § 3).
\end{enumerate}

\exercisesection{Smooth algebras}

\begin{enumerate}
\def\labelenumi{\arabic{enumi})}
\tightlist
\item
  Let \(k\) a ring, A a linearly topologized formally smooth \(k\) -linearly topologized algebra. One says that A is a \(k\) -algebra formally étale (resp. formally net) if it satisfies the following condition: whatever the \(k\) -algebra C, endowed with the discrete topology, and the ideal N of square zero of C, every continuous homomorphism of \(k\) -algebras A → C/N admits a unique (resp. at most one) lifting A → C. For A to be formally étale over \(k\) it is necessary and sufficient that it be formally smooth and formally unramified. Every quotient of \(k\) is formally unramified.
\end{enumerate}

\begin{enumerate}
\def\labelenumi{\alph{enumi})}
\item
  Prove the analogues of the statements of Props. 3 and 4 for formally smooth or formally étale algebras.
\item
  Let J be an ideal of \(\Lambda\) such that the topology of A is the J-adic topology; denote by \(\widehat{\Omega}_{k}(\Lambda)\) the separated completion of the A-module \(\Omega_{k}(A)\) (for the J-adic topology). For the \(k\) -topological algebra A to be formally smooth, it is necessary and sufficient that \(\widehat{\Omega}_{k}(A)\) be zero (observe that each of these properties is equivalent to the fact that every \(k\) -derivation of A into an A-module annihilated by a power of J is zero).
\end{enumerate}

\begin{enumerate}
\def\labelenumi{\arabic{enumi})}
\setcounter{enumi}{1}
\item
  Let A be a Noetherian ring, J an ideal of A, and \(\widehat{\mathbf{A}}\) the separated completion of A for the J-adic topology. If the \(\Lambda\) -algebra \(\widehat{\mathbf{A}}\) is formally smooth (for the discrete topology), it is formally étale (observe that one has \(\Omega_{\mathbf{A}}(\widehat{\mathbf{A}}) = J^{n}\Omega_{\mathbf{A}}(\widehat{\mathbf{A}})\) for every integer \(n\) ).
\item
  Let K be a finitely generated extension of a field \(k\) ; prove that the following conditions are equivalent:
\end{enumerate}

\begin{enumerate}
\def\labelenumi{(\roman{enumi})}
\item
  the k-algebra K is formally smooth;
\item
  the \(k\) -algebra K is formally étale;
\item
  K is a finite separable extension of k.
\end{enumerate}

\begin{enumerate}
\def\labelenumi{\arabic{enumi})}
\setcounter{enumi}{3}
\tightlist
\item
  Let \(k\) be a Noetherian ring, A a finitely generated \(k\) -algebra. One says that \(\Lambda\) is an \(k\) -étale algebra (resp. smooth) if it is formally étale (resp. formally smooth) when equipped with the discrete topology.
\end{enumerate}

\begin{enumerate}
\def\labelenumi{\alph{enumi})}
\tightlist
\item
  Prove that the following conditions are equivalent:
\end{enumerate}

\begin{enumerate}
\def\labelenumi{(\roman{enumi})}
\item
  the k-algebra A is smooth;
\item
  \(\Omega_k(A) = 0\)
\item
  the \(\mathrm{A}\otimes_{k}\mathrm{A}\) -module A is projective.
\end{enumerate}

(Let I be the kernel of the canonical homomorphism \(\mathrm{A} \otimes_k \mathrm{A} \to \mathrm{A}\) ; show that (ii) and (iii) are both equivalent to

\begin{enumerate}
\def\labelenumi{(\roman{enumi})}
\setcounter{enumi}{3}
\tightlist
\item
  \(I_{q} = 0\) for every prime ideal q of A \(\otimes_{k}\) A containing I.)
\end{enumerate}

\begin{enumerate}
\def\labelenumi{\alph{enumi})}
\setcounter{enumi}{1}
\item
  Let \(f \in k[T]\) be a polynomial such that f and \(f'\) generate the unit ideal of k{[}T{]}. Prove that the k-algebra \(k[T]/(f)\) is étale (cf.~no. 3, Example 3).
\item
  If k is a field, the smooth k-algebras are the étale k-algebras, and this notion coincides with that introduced in A, V, p.~28 (the case of algebras of finite degree over k follows from loc. cit., p.~32, th. 3; if A is a smooth k-algebra, consider a maximal ideal m of A and apply the preceding case to A/m \(^{n}\) ). The smooth k-algebras are therefore the finite products of finite separable extensions of k (A, V, p.~34, th. 4).
\item
  For the k-algebra A to be smooth, it is necessary and sufficient that for every prime ideal p of k the \(\kappa(\mathfrak{p})\) -algebra \(\kappa(\mathfrak{p})\otimes_{k}A\) be étale.
\end{enumerate}

\begin{enumerate}
\def\labelenumi{\arabic{enumi})}
\setcounter{enumi}{4}
\tightlist
\item
  Let A be a complete Noetherian local ring and B an essentially finitely generated A-algebra.
\end{enumerate}

\begin{enumerate}
\def\labelenumi{\alph{enumi})}
\item
  Show that the set of prime ideals p of A such that the ring \(\Lambda_{\mathfrak{p}}\) is regular is open in \(\operatorname{Spec}(\mathbf{A})\) (one will use Exercise 16 of VIII, § 5 to reduce to the case where A is integral, then IX, § 2, no. 5, th. 3 and Exercise 17 of VIII, § 5).
\item
  Suppose that A contains a field \(^{1}\) . Show that the set of prime ideals q of B such that the ring \(B_{q}\) is regular is open in Spec(B) (one will use § 7, no. 9, cor. 2 of th. 3 and IX, § 3, no. 3, th. 2).
\end{enumerate}

\begin{enumerate}
\def\labelenumi{\arabic{enumi})}
\setcounter{enumi}{5}
\item
  Let \(k\) a ring, A a linearly topologized formally smooth \(k\) -algebra. Prove that for every open ideal J of A, the A/J-module \((\mathrm{A} / \mathrm{J}) \otimes_{\mathrm{A}} \Omega_k(\mathrm{A})\) is projective (prove, using the example of no. 1, that for every surjective homomorphism of A/J-modules \(\mathrm{M} \to \mathrm{M}''\) , every \(k\) -derivation of A into \(\mathrm{M}''\) lifts to M).
\item
  Let \(k\) a ring, \(\rho : A \to B\) a homomorphism of \(k\) -algebras, \(J\) an ideal of \(B\) . One says that \(B\) , equipped with the \(J\) -adic topology, is formally smooth over \(A\) relative to \(k\) if for every \(A\) -algebra \(C\) (which we endow with the discrete topology) and every ideal of square zero \(N\) of \(C\) , every continuous homomorphism of \(A\) -algebras \(B \to C/N\) that lifts to a continuous homomorphism of \(k\) -algebras from \(B\) in \(C\) also lifts to a continuous homomorphism of \(A\) -algebras from \(B\) in \(C\) .
\end{enumerate}

\begin{enumerate}
\def\labelenumi{\alph{enumi})}
\tightlist
\item
  Show that the following conditions are equivalent:
\end{enumerate}

\begin{enumerate}
\def\labelenumi{(\roman{enumi})}
\item
  B is formally smooth over A relative to k for the J-adic topology.
\item
  for all \(k\) -derivation D of \(\Lambda\) In a B-module E annihilated by a power of J, there exists a \(k\) -derivation \(\tilde{D}: B \to E\) such that \(\tilde{D} \circ \rho = D\) ;
\item
  for every integer \(n \geqslant 0\) , the canonical map \(u_n: \Omega_k(A) \otimes_A (B/J^n) \longrightarrow \Omega_k(B) \otimes_B (B/J^n)\) is injective and its image is a direct summand.
\end{enumerate}

\begin{enumerate}
\def\labelenumi{\alph{enumi})}
\setcounter{enumi}{1}
\tightlist
\item
  We assume B is formally smooth over k for the J-adic topology. Prove that for B to be formally smooth over A, it is necessary and sufficient that the canonical homomorphism \(u_{1}:\Omega_{k}(A)\otimes_{A}(B/J)\to\Omega_{k}(B)\otimes_{B}(B/J)\) be injective and that its image be a direct summand (use a) and exerc. 6 to prove that a retraction of \(u_{1}\) lifts to a retraction of \(u_{n}\) for every n).
\end{enumerate}

¶ 8) Let k be a field and T a finite family of indeterminates. One equips the k-algebra \(k[[T]]\) with the discrete topology.

\begin{enumerate}
\def\labelenumi{\alph{enumi})}
\item
  Prove that \(k[[T]]\) is formally smooth over \(k[T]\) with respect to k (exerc. 7).
\item
  If \(k[[\mathbf{T}]]\) is formally smooth over \(k\) , prove that the characteristic \(p\) of \(k\) is strictly positive and that \(k\) is a finite extension of \(k^p\) (deduce from exerc. 2 that \(\Omega_{k|\mathbf{T}|}(k[[\mathbf{T}]])\) is zero, and consequently that \(\Omega_{k(\mathbf{T})}(k((\mathbf{T})))\) is zero; then prove, using prop. 6 of A, V, p.~99, that the subextension of \(k((\mathbf{T}))\) consisting of series whose coefficients generate a finitely generated extension of \(k^p\) is equal to \(k((\mathbf{T})))\) .
\end{enumerate}

\begin{enumerate}
\def\labelenumi{\arabic{enumi})}
\setcounter{enumi}{8}
\tightlist
\item
  Let \(k\) a field of characteristic \(p > 0\) , containing a decreasing filtered family \((k_{\lambda})_{\lambda \in \Lambda}\) of subfields satisfying \(\bigcap_{\lambda \in \Lambda} k_{\lambda} = k^{p}\) .
\end{enumerate}

\begin{enumerate}
\def\labelenumi{\alph{enumi})}
\item
  Prove that the intersection of the kernels of the canonical homomorphisms from \(\Omega(k)\) in \(\Omega_{k_{\lambda}}(k)\) , for \(\lambda\) traversing \(\Lambda\) is reduced to (0) (consider a \(p\) -basis of \(k\) ).
\item
  Let A be a regular local k-algebra, endowed with the topology \(m_{A}\) -adic. Prove that for A to be formally smooth over k, it is necessary and sufficient that \(\Lambda\) let be formally smooth over k relative to \(k_{\lambda}\) , for all \(\lambda \in \Lambda\) (using a), Prop. 8, and Exerc. 7, a)).
\end{enumerate}

¶ 10) Let A be a complete Noetherian local ring, K its field of fractions. Let p be a prime ideal of A; denote by B the local ring \(A_{p}\) and \(\widehat{B}\) its completion.

\begin{enumerate}
\def\labelenumi{\alph{enumi})}
\item
  Assume A is regular and K has characteristic zero. Show that the K-algebra \(K \otimes_{B} \widehat{B}\) is absolutely regular.
\item
  Assume A is regular and K has characteristic p. \textgreater{} 0, so that A is identified with a ring of formal power series. \(k[[T_{1}, \ldots, T_{n}]]\) (IX, § 3, no. 3, th. 2). We denote by \((k_{\lambda})_{\lambda \subset \Lambda}\) the family of subfields \(k_{\lambda}\) of k containing \(k^{p}\) and such that k is a finite-degree extension of \(k_{\lambda}\) . Let \(\lambda \in \Lambda\) ; we denote \(A_{\lambda}\) the ring \(k_{\lambda}[[T_{1}^{p}, \ldots, T_{n}^{p}]]\) , \(B_{\lambda}\) the local ring of \(A_{\lambda}\) into the prime ideal \(p \cap A_{\lambda}\) , and by \(K_{\lambda}\) the field of fractions of \(\Lambda_{\lambda}\) and \(B_{\lambda}\) Prove that the A-algebra B (endowed with the discrete topology) is formally smooth relative to \(B_{\lambda}\) (one may first show that \(\widehat{B}\) is isomorphic to \(B \otimes_{B_{\lambda}} \widehat{B_{\lambda}}\) ).
\item
  Under the hypotheses of b), prove that every local ring of \(K \otimes_{B} \widehat{B}\) is of the form \((\widehat{B})_{\mathbf{r}}\) , with \(\mathfrak{r} \in \operatorname{Spec}(\widehat{B})\) and \(\mathfrak{r} \cap B = \{0\}\) , and that it is formally smooth over K relative to \(K_{\lambda}\) Deduce from this that the K-algebra \(K \otimes_{B} \widehat{B}\) is absolutely regular (prove that the family \(\{K_{\lambda}\}_{\lambda \in \Lambda}\) of subfields of K is downward filtering and that \(\bigcap_{\lambda \in \Lambda} K_{\lambda} = k^{p}((T_{1}, \ldots, T_{n}))\) ; then apply exercise 9).
\item
  We return to the general case. Prove that A is a Grothendieck ring (§ 6, exerc. 5); if \(\mathfrak{q} \in \operatorname{Spec}(B)\) , we need to show that the \(\kappa(\mathfrak{q})\) -algebra \(\kappa(\mathfrak{q}) \otimes_{\mathrm{B}} \widehat{\mathrm{B}}\) is absolutely regular; replacing A by \(\mathrm{A}/(\mathrm{A} \cap \mathfrak{q})\) one may assume \(\Lambda\) integral domain and \(\mathfrak{q} = \{0\}\) ; we will use IX, § 2, n° 3, th. 3 and § 3, n° 3, th. 2 to reduce to the case where A is a complete regular local ring.
\end{enumerate}

\begin{enumerate}
\def\labelenumi{\arabic{enumi})}
\setcounter{enumi}{10}
\item
  Let A be a Noetherian ring; suppose that for every maximal ideal m of A, the completion of the local ring \(A_{m}\) is an \(A_{m}\) absolutely regular algebra. Show that A is a Grothendieck ring (§ 6, exerc. 5). (Use loc. cit. and the preceding exercise to show that for every maximal ideal m of A, the ring \(A_{m}\) is a Grothendieck ring.)
\item
  Let A be a Noetherian local ring. For A to be a Grothendieck ring, it is necessary and sufficient that, for every finite integral A-algebra B, every maximal ideal m of B, and every prime ideal q of the completed ring \(\widehat{B_{m}}\) satisfying \(B_{m} \cap q = (0)\) , the local ring \((\widehat{B_{m}})_{q}\) let be regular (to show that this condition is necessary, it suffices, by the previous exercise, to show that, for every prime ideal p of A, and every finite extension K of \(\kappa(\mathfrak{p})\) , the ring \(K \otimes_{A} \widehat{A}\) is regular; one may introduce a finite integral A-algebra B with fraction field \(K\) to note \(\mathfrak{m}_{1},\ldots,\mathfrak{m}_{r}\) its maximal ideals, and interpret the local rings of \(K\otimes_{A}\widehat{A}\) as local rings of one of the \(\widehat{B_{\mathfrak{m}_{i}}}\) .
\item
  Let \(k\) a field and A a \(k\) -algebra essentially of finite type. Show that A is a Grothendieck ring (reduce to the case where A is a local ring of a polynomial ring over \(k\) in a finite number of indeterminates; using the preceding exercise, reduce next to showing that, for every prime ideal \(\mathfrak{r}\) of \(\mathrm{A}[\mathrm{T}_1,\ldots ,\mathrm{T}_n]\) , every local ring C of \(\mathrm{A}[\mathrm{T}_1,\ldots ,\mathrm{T}_n]\) at a maximal ideal containing \(\mathfrak{r}\) , and every ideal \(\mathfrak{q}\) of the completed ring \(\widehat{\mathrm{C}}\) such that \(\mathrm{C}\cap \mathfrak{q} = \mathfrak{r}\mathrm{C}\) , the ring \((\widehat{\mathrm{C}})_q / \mathfrak{r}(\widehat{\mathrm{C}})_q\) is regular. For this, one may use th. 3 of no. 9).
\end{enumerate}

¶ 14) Let A be a local Grothendieck ring. Show that the set of prime ideals p of A such that the ring \(A_{p}\) is regular is open in \(\operatorname{Spec}(A)\) (let \(\widehat{A}\) the completion of A; we will use exercise 5, exercise 16 of VIII, § 5, and the fact that the canonical mapping \(\operatorname{Spec}(\widehat{A}) \to \operatorname{Spec}(A)\) is open).

¶ 15) Let \(k\) a perfect field of characteristic \(p > 0\) , \((\mathrm{T}_{n})_{n \in \mathbf{N}}\) a family of indeterminates, and K the field \(k((\mathrm{T}_{n})_{n \in \mathbf{N}})\) . Let \(\Lambda\) the subring of the ring of formal power series \(\mathrm{K}[[\mathrm{X},\mathrm{Y}]]\) generated by \(\mathrm{K}^{p}[[\mathrm{X},\mathrm{Y}]]\) and K. The ring A is a regular local Noetherian ring of dimension 2, and its completion is \(\mathrm{K}[[\mathrm{X},\mathrm{Y}]]\) For every integer \(n\) , we denote \(p_{n}\) the element \(\mathrm{X} - \mathrm{T}_{n}\mathrm{Y}\) of A, \(q_{n}\) the product \(p_{0} \ldots p_{n}\) , and we set \(c = \sum_{n \in \mathbf{N}} q_{n}\mathrm{T}_{n}\) We denote by B the sub-A-algebra dc \(\mathrm{K}[[\mathrm{X},\mathrm{Y}]]\) generated by \(c\) .

\begin{enumerate}
\def\labelenumi{\alph{enumi})}
\item
  Let \(n\) an integer \(> 0\) and \(\mathfrak{p}\) a prime ideal of \(B\) of height 1 that contains \(p_n\) Show that the ring \(B_{\mathfrak{p}}\) is not normal (for every integer \(m\) , we will set \(c_m = (c - \sum_{i=0}^{m-1} q_i T_i)/q_m\) and we will show that \(B_{\mathfrak{p}} = A_{A p_n}[c_{n-1}]\) , so that \(c_n\) is not in \(B_{\mathfrak{p}}\) , whereas \(c_n^p\) belongs to \(A\) ).
\item
  Deduce that the set of prime ideals p of B such that the ring \(B_{p}\) is regular contains no nonempty open subset of Spec(B) (note that B is a faithfully flat A-module).
\end{enumerate}

\exercisesection{Duality of modules of finite length}

\begin{enumerate}
\def\labelenumi{\arabic{enumi})}
\tightlist
\item
  Let A be a Noetherian local ring, and M an A-module. Show that for M to be Artinian, it is necessary and sufficient that its socle S be of finite dimension over \(\kappa_{\Lambda}\) and that M be an essential extension of S (that is, by (A, II, p.~185, exerc. 15), every nonzero submodule of M contains a nonzero element of S).
\end{enumerate}

¶ 2) Let A be a complete local Noetherian ring, and let M be a Λ-module. Show that the following conditions are equivalent:

\begin{enumerate}
\def\labelenumi{(\roman{enumi})}
\item
  the canonical homomorphism \(\alpha_{\mathrm{M}}: \mathrm{M} \to \mathrm{D}(\mathrm{D}(\mathrm{M}))\) is bijective;
\item
  there exists an exact sequence \(0 \to M' \to M \to M'' \to 0\) , where \(M'\) is of finite type and \(M''\) Artinian.
\end{enumerate}

(If M satisfies (i), show that its socle is of finite dimension over \(\kappa_{A}\) , and that the same is true of every quotient of M; construct a finitely generated submodule N of M such that M/mN is an essential extension of its socle, and use exerc. 1).

\begin{enumerate}
\def\labelenumi{\arabic{enumi})}
\setcounter{enumi}{2}
\tightlist
\item
  Let \(k\) a field, S the ring \(k[\mathrm{T}_{1},\ldots ,\mathrm{T}_{d}]\) , \(\mathfrak{m}\) the ideal \((\mathrm{T}_1,\dots ,\mathrm{T}_d)\) . One denotes by \(\mathrm{S^{*gr}}\) the graded dual of S (no. 6), and \((\mathbf{T}^{-\alpha})_{\alpha \in \mathbf{N}^d}\) the dual basis of \((\mathbf{T}^{\alpha})_{\alpha \in \mathbf{N}^d}\) (denoted \((u_{\alpha})_{\alpha \in \mathbf{N}^d}\) in loc. cit.).
\end{enumerate}

\begin{enumerate}
\def\labelenumi{\alph{enumi})}
\item
  To every finitely generated sub-S-module M of S*gr one associates the ideal Ann(M) of S. Show that this defines a bijective correspondence between the finitely generated submodules of S*gr and the ideals a of S contained in m and such that S/a is of finite length; the inverse bijection associates to the ideal a the submodule M of S*gr consisting of the elements annihilated by a (observe that M is a Matlis S/a-module).
\item
  For S/a to be a Gorenstein ring, it is necessary and sufficient that the S-module M be cyclic.
\item
  Let \(f\) the element \(\sum \mathrm{T}_i^{-2}\) of \(\mathrm{S}^{*\mathrm{gr}}\) ; prove that the ideal \(\mathfrak{a}\) corresponding to the submodule \(A f\) of \(\mathrm{S}^{*\mathrm{gr}}\) is generated by the quadratic forms \(\mathrm{T}_i\mathrm{T}_j\) for \(i < j\) and \(\mathrm{T}_i^2 - \mathrm{T}_1^2\) for \(i > 1\) If \(d \geqslant 3\) the ideal \(\mathfrak{a}\) is not completely secant; the Gorenstein ring \(A / \mathfrak{a}\) is not a complete intersection ring (cf. § 5, exerc. 5).
\end{enumerate}

\begin{enumerate}
\def\labelenumi{\arabic{enumi})}
\setcounter{enumi}{3}
\item
  Let A be a Noetherian ring, M a finitely generated A-module. For every \(\mathfrak{p} \in \operatorname{Spec}(A)\) , we set \(\boldsymbol{\mu}^{i}(\mathfrak{p}, M) = \dim_{\kappa(\mathfrak{p})} \operatorname{Ext}_{A_{\mathfrak{p}}}^{i}(\kappa(\mathfrak{p}), M_{\mathfrak{p}})\) ; we denote \((I(\mathfrak{p}), e_{\mathfrak{p}})\) an injective envelope of the A-module \(A/\mathfrak{p}\) . Let I be a minimal injective resolution of M (which means that \(I^{n}\) is an injective envelope of \(Z^{n}(I)\) for every \(n \geqslant 0\) , cf. A, X, p. 182, exerc. 13). Prove that the A-module \(I^{p}\) is isomorphic to \(\bigoplus_{\mathfrak{p}} I(\mathfrak{p})^{\mu^{p}(\mathfrak{p}, M)}\) .
\item
  Let A be a Noetherian local ring, I a Matlis A-module. For every A-module M one denotes by D(M) the Matlis dual \(\mathrm{Hom}_{\mathrm{A}}(\mathrm{M},\mathrm{I})\) . Recall that one denotes by \(\mathrm{dpl}_{\mathrm{A}}(\mathrm{M})\) the infimum of the lengths of flat resolutions of M (A, X, p. 202, exerc. 7). Prove the equalities \(\mathrm{di}_{\mathrm{A}}(\mathrm{D}(\mathrm{M})) = \mathrm{dpl}_{\mathrm{A}}(\mathrm{M})\) , \(\mathrm{di}_{\mathrm{A}}(\mathrm{M}) = \mathrm{dpl}_{\mathrm{A}}(\mathrm{D}(\mathrm{M}))\) .
\item
  Let A be a Noetherian ring, I and J injective A-modules. Prove that the A-module \(\mathrm{Hom}_{\mathrm{A}}(\mathrm{I},\mathrm{J})\) is flat (cf. prop. 6, b)).
\item
  Let A and B be rings not necessarily commutative, J a (A, B)-bimodule, P a flat left B-module. Assume the ring A is left Noetherian. Prove that if J is an injective A-module, then J ⊗B P is likewise.
\end{enumerate}

¶ 8) Let A be a Noetherian local ring.

\begin{enumerate}
\def\labelenumi{\alph{enumi})}
\tightlist
\item
  Let M and N be A-modules such that \(\mathrm{dpl}_{\mathrm{A}}(\mathrm{M}) < +\infty\) and \(\operatorname{Tor}_{i}^{\mathrm{A}}(\mathrm{M}, \mathrm{N}) = 0\) for i \textgreater{} 0. Define for every integer n a canonical isomorphism
\end{enumerate}

\[
\operatorname{Ext} _ {\mathrm{A}} ^ {n} \left(\kappa_ {\mathrm{A}}, \mathrm{M} \otimes_ {\mathrm{A}} \mathrm{N}\right) \longrightarrow \bigoplus_ {p \geqslant n} \left(\operatorname{Tor} _ {p - n} ^ {\mathrm{A}} \left(\kappa_ {\Lambda}, \mathrm{M}\right) \otimes_ {k} \operatorname{Ext} _ {\mathrm{A}} ^ {p} \left(\kappa_ {\mathrm{A}}, \mathrm{N}\right)\right),
\]

and prove that one has \(\mathrm{di}_{\mathrm{A}}(\mathrm{M}\otimes_{\mathrm{A}}\mathrm{N})\leqslant \mathrm{di}_{\mathrm{A}}(\mathrm{N})\) (adapt the proof of exerc. 6 of § 3 with the aid of exerc. 7).

\begin{enumerate}
\def\labelenumi{\alph{enumi})}
\setcounter{enumi}{1}
\tightlist
\item
  Let N, P be A-modules such that \(\mathrm{di}_{\mathrm{A}}(\mathrm{P}) < +\infty\) and \(\mathrm{Ext}_{\mathrm{A}}^{i}(\mathrm{N},\mathrm{P}) = 0\) for i \textgreater{} 0. Define for every integer n a canonical isomorphism
\end{enumerate}

\[
\operatorname{Tor} _ {n} ^ {\mathrm{A}} \left(\kappa_ {\mathrm{A}}, \operatorname{Hom} _ {\mathrm{A}} (\mathrm{N}, \mathrm{P})\right) \longrightarrow \bigoplus_ {p \geqslant n} \operatorname{Hom} _ {\kappa_ {\mathrm{A}}} \left(\operatorname{Ext} _ {\mathrm{A}} ^ {p - n} \left(\kappa_ {\mathrm{A}}, \mathrm{P}\right), \operatorname{Ext} _ {\mathrm{A}} ^ {p} \left(\kappa_ {\mathrm{A}}, \mathrm{N}\right)\right)
\]

(apply Matlis duality).

\begin{enumerate}
\def\labelenumi{\alph{enumi})}
\setcounter{enumi}{2}
\item
  Prove that if there exists an A-module M whose injective dimension and flat dimension are finite, then A is a Gorenstein ring (take N = A in a)).
\item
  Conversely, if A is a Gorenstein ring, prove that there is an identity between modules of finite injective dimension and modules of finite flat dimension (prove with the aid of a) that a module of finite flat dimension is of finite injective dimension, and obtain the converse by Matlis duality).
\item
  Deduce from d) that the following conditions are equivalent:
\end{enumerate}

\begin{enumerate}
\def\labelenumi{(\roman{enumi})}
\item
  A is a Gorenstein ring;
\item
  the finitely generated A-modules of finite injective dimension coincide with the finitely generated A-modules of finite projective dimension;
\item
  there exists a finitely generated A-module whose injective dimension and projective dimension are finite.
\end{enumerate}

\begin{enumerate}
\def\labelenumi{\arabic{enumi})}
\setcounter{enumi}{8}
\tightlist
\item
  Let A be a Noetherian local ring satisfying condition (GM) (§ 3, exerc. 13).
\end{enumerate}

\begin{enumerate}
\def\labelenumi{\alph{enumi})}
\item
  Let I be a complex of injective A-modules, of length \(\ell\) , whose homology is nonzero and of finite type. Prove the inequality \(\ell \geqslant \dim(A)\) (reduce to loc. cit. a) by Matlis duality).
\item
  If there exists a finitely generated A-module M of finite injective dimension, then A is a Macaulay ring (apply a) to an injective resolution of M).
\item
  For A to be a Gorenstein ring, it is necessary and sufficient that it possess a cyclic module of finite injective dimension (use Exercise 8, b)).
\end{enumerate}

\begin{enumerate}
\def\labelenumi{\arabic{enumi})}
\setcounter{enumi}{9}
\tightlist
\item
  Let A be a Noetherian ring, and let M, P, T be A-modules; suppose that M is finitely generated, T has finite injective dimension, and that the modules \(\operatorname{Ext}_{\mathrm{A}}^{i}(\mathrm{M},\mathrm{P})\) and \(\operatorname{Ext}_{\mathrm{A}}^{i}(\mathrm{P},\mathrm{T})\) are zero for i \textgreater{} 0. Prove that the canonical homomorphism \(\mathrm{M}\otimes_{\mathrm{A}}\mathrm{Hom}_{\mathrm{A}}(\mathrm{P},\mathrm{T})\to\mathrm{Hom}_{\mathrm{A}}(\mathrm{Hom}_{\mathrm{A}}(\mathrm{M},\mathrm{P}),\mathrm{T})\) which maps \(m\otimes h\) onto the homomorphism \(u\mapsto h(u(m))\) is an isomorphism (argue as in the proof of Prop. 6, replacing J by an injective resolution of T).
\end{enumerate}

\exercisesection{Dualizing modules}

\begin{enumerate}
\def\labelenumi{\arabic{enumi})}
\tightlist
\item
  Let A be a local Macaulay ring.
\end{enumerate}

\begin{enumerate}
\def\labelenumi{\alph{enumi})}
\item
  We assume that A admits a dualizing module. \(\Omega\) ; we endow the A-module \(\mathrm{A} \oplus \Omega\) of the A-algebra structure for which \(\Omega\) is an ideal with square zero. Prove that \(\mathrm{A} \oplus \Omega\) is a Gorenstein ring (using a maximal regular sequence, reduce to the case where A is Artinian; observe that the socle of \(\mathrm{A} \oplus \Omega\) is identified with that of the A-module \(\Omega\) ).
\item
  Deduce that for A to admit a dualizing module, it is necessary and sufficient that there exist a Gorenstein ring B and a surjective homomorphism from B to A.
\end{enumerate}

\begin{enumerate}
\def\labelenumi{\arabic{enumi})}
\setcounter{enumi}{1}
\tightlist
\item
  Let A be a regular local ring, I an ideal of A, B the ring A/I; suppose that B is a Macaulay ring. Let (L,p) be a minimal projective resolution of the A-module B; it has length \(c = \text{ht}(I)\) .
\end{enumerate}

\begin{enumerate}
\def\labelenumi{\alph{enumi})}
\item
  Let \(\mathrm{L}^*\) the complex \(\operatorname{Homgr}_{\mathrm{A}}(\mathrm{L},\mathrm{A})\) ; prove that we have \(\mathrm{H}^{i}(\mathrm{L}^{*}) = 0\) for \(i\neq c\) and that the B-module \(\mathrm{H}^c (\mathrm{L}^*)\) is dualizing.
\item
  Prove that the following conditions are equivalent:
\end{enumerate}

\begin{enumerate}
\def\labelenumi{(\roman{enumi})}
\item
  A is a Gorenstein ring;
\item
  the complexes L and \(L^{*}(-d)\) are isomorphic;
\item
  \(L_{c}\) is free of rank one.
\end{enumerate}

\begin{enumerate}
\def\labelenumi{\arabic{enumi})}
\setcounter{enumi}{2}
\tightlist
\item
  Let A be a ring admitting a dualizing module \(\Omega\) .
\end{enumerate}

\begin{enumerate}
\def\labelenumi{\alph{enumi})}
\item
  For \(\Omega\) is isomorphic to an ideal I of A, it is necessary and sufficient that \(A_{p}\) let R be a Gorenstein ring for every minimal prime ideal \(p\) of A (observe that \(\Omega\) is isomorphic to a submodule of \(\bigoplus \Omega_{p}\) ).
\item
  When these conditions are satisfied, prove that I is equal to A or has height 1, and that A/I is a Gorenstein ring (compute the depth of A/I, then the A/I-module \(\mathrm{Ext}_{\mathrm{A}}^{1}(\mathrm{A / I,I})\) ).
\item
  Prove that a factorial ring admitting a dualizing module is a Gorenstein ring.
\end{enumerate}

\begin{enumerate}
\def\labelenumi{\arabic{enumi})}
\setcounter{enumi}{3}
\item
  Let A be a local Macaulay ring of dimension \(d\) ; let \(\Omega\) a finitely generated A-module, of finite injective dimension, of depth \(d\) , such that the canonical homomorphism \(\mathrm{A} \to \mathrm{End}_{\mathrm{A}}(\Omega)\) is bijective. Prove that the \(\Lambda\) -module \(\Omega\) is dualizing (argue by induction on \(d\) , deducing from exerc. 7 of § 3 that for every element \(\Omega\) -regular \(x\) of \(\mathfrak{m}_{\mathrm{A}}\) , the canonical homomorphism \(\mathrm{A}/x\mathrm{A} \to \mathrm{End}_{\mathrm{A}/x\mathrm{A}}(\Omega/x\Omega)\) is bijective).
\item
  Let A be a Noetherian local ring admitting a dualizing module \(\Omega\) . Let T be a finitely generated A-module of finite injective dimension.
\end{enumerate}

\begin{enumerate}
\def\labelenumi{\alph{enumi})}
\item
  Prove that \(\mathrm{Ext}_{\mathrm{A}}^{i}(\Omega, \mathrm{T})\) is zero for \(i > 0\) (use exerc. 7 of § 3).
\item
  Prove that the canonical homomorphism \(\Omega\otimes_{\Lambda}\mathrm{Hom}_{\mathbb{A}}(\Omega,\mathrm{T})\to\mathrm{T}\) is bijective (apply a) and exerc. 10 of § 8).
\item
  Prove using exercise 8, b) of § 8 the equalities dp \(_{A}\) Hom \(_{A}\) (Ω,T)) = dim(A) - prof(T) and prof(Hom \(_{A}\) (Ω,T)) = prof(T). Deduce that if prof(T) = dim(A), then the A-module T is a direct sum of modules isomorphic to Ω.
\end{enumerate}

\(d)\) We set \(c = \dim(A) - \operatorname{prof}(T)\) Prove that there exist integers \(n_0, \ldots, n_c\) and an exact sequence

\[
0 \rightarrow \Omega^ {n _ {c}} \rightarrow \dots \rightarrow \Omega^ {n _ {0}} \rightarrow \mathrm{T} \rightarrow 0
\]

(argue by induction on c, constructing with the aid of b) a surjective homomorphism \(\Omega^{n_{0}} \to T\) ).

¶ 6) Let A be a local noetherian ring admitting a dualizing module Ω; let M be a finitely generated A-module of finite projective dimension.

\begin{enumerate}
\def\labelenumi{\alph{enumi})}
\item
  Prove that one has \(\operatorname{Tor}_{i}^{\mathrm{A}}(\Omega, \mathrm{M}) = 0\) for \(i > 0\) (Argue by induction on \(\dim(\mathrm{A})\) , considering an exact sequence \(0 \to \mathrm{N} \to \mathrm{L} \to \mathrm{M} \to 0\) where \(\mathrm{L}\) is free of finite type. If \(x\) is a non-zero-divisor element of \(\mathfrak{m}_{\mathrm{A}}\) , deduce from the induction hypothesis that \(\operatorname{Tor}_{i}^{\mathrm{A}}(\Omega, \mathrm{N})\) is zero for \(i > 0\) , which implies \(\operatorname{Tor}_{i}^{\mathrm{A}}(\Omega, \mathrm{M}) = 0\) for \(i > 1\) , and that the homothety of ratio \(x\) is injective in \(\Omega \otimes_{\mathrm{A}} \mathrm{N}\) , which implies that every prime ideal \(\mathfrak{p}\) associated with \(\operatorname{Tor}_{1}^{\mathrm{A}}(\Omega, \mathrm{M})\) is a minimal ideal of \(\operatorname{Spec}(\mathrm{A})\) ; observe that at such a \(\mathfrak{p}\) the \(\mathrm{A}_{\mathfrak{p}}\) -module \(\mathrm{M}_{\mathfrak{p}}\) is free).
\item
  Prove that the A-module \(\Omega\otimes_{\mathrm{A}}\mathrm{M}\) is of finite injective dimension.
\item
  Prove that the map \(\theta: \mathrm{M} \to \mathrm{Hom}_{\Lambda}(\Omega, \Omega \otimes_{\Lambda} \mathrm{M})\) defined by \(\theta(m)(w) = w \otimes m\) for \(m \in \mathrm{M}\) , \(w \in \Omega\) is an isomorphism (reduce with the aid of a projective resolution to the case where M is free).
\end{enumerate}

\(d\) ) The maps \([\mathrm{T}] \mapsto [\mathrm{Hom}_{\mathrm{A}}(\Omega, \mathrm{T})]\) and \([\mathrm{M}] \mapsto [\Omega \otimes_{\mathrm{A}} \mathrm{M}]\) define mutually inverse bijections between the set of isomorphism classes of finitely generated A-modules of finite injective dimension and the set of isomorphism classes of finitely generated A-modules of finite projective dimension, * and quasi-inverse equivalences of categories between the corresponding categories*.

¶ 7) Let A be a macaulay local ring, and M a finitely generated macaulay A-module. We denote by \(r(M)\) (resp. \(t(M)\) ) the dimension of the \(\kappa_{A}\) -vector space \(\kappa_{A} \otimes_{A} M\) (resp. \(\operatorname{Ext}_{A}^{p}(\kappa_{A}, M)\) , with \(p = \operatorname{prof}(M)\) ), cf. exercise 16 of § 1.

\begin{enumerate}
\def\labelenumi{\alph{enumi})}
\item
  Suppose that A possesses a dualizing module \(\Omega\) ; set \(M' = \operatorname{Ext}_A^c(M, \Omega)\) , with \(c = \dim(A) - \dim_A(M)\) . Prove the equalities \(r(M') = t(M)\) and \(t(M') = r(M)\) (reduce with the aid of prop. 7 of § 3 to the case \(c = 0\) , then by passing to the quotient by a maximal secant sequence to the case where A is artinian, and use prop. 6 of § 8).
\item
  Hence we have \(t(\mathrm{M}_{\mathfrak{p}}) \leqslant t(\mathrm{M})\) for every \(\mathfrak{p} \in \operatorname{Spec}(\mathrm{A})\) (reduce with the aid of exercises 2 of § 2 and 16 of § 1 to the case where A is complete, and apply a)).
\item
  For the module M to be dualizing, it is necessary and sufficient that it be faithful and that one have \(t(\mathrm{M}) = 1\) (reduce to the case where A is complete; apply a) and the corollary of th. 1).
\end{enumerate}

\begin{enumerate}
\def\labelenumi{\arabic{enumi})}
\setcounter{enumi}{7}
\tightlist
\item
  Let A be a local Gorenstein ring, a a non-zero ideal of height zero, and b the annihilator of a.
\end{enumerate}

\begin{enumerate}
\def\labelenumi{\alph{enumi})}
\item
  Prove that the ideal b is of height zero, and that the A-module A/b has no embedded associated prime ideals (observe that every prime ideal associated to A/b is associated to A).
\item
  Suppose that the A-module A/α has no embedded associated prime ideals; prove that α is the annihilator of β (reduce to the case where dim(A) = 0, and use Matlis duality).
\item
  For A/a to be a Macaulay ring, it is necessary and sufficient that the same be true of A/b (use the corollary of th. 1 and that of prop. 3). If these conditions are satisfied, the A/a-module b is dualizing.
\end{enumerate}

\begin{enumerate}
\def\labelenumi{\arabic{enumi})}
\setcounter{enumi}{8}
\tightlist
\item
  Let \(\Gamma\) a submonoid of the additive monoid \(\mathbf{N}\) such that \(\mathbf{N} - \Gamma\) be finite; we denote by \(c\) the smallest integer such that \(\Gamma\) contains all the integers \(\geqslant c\) . Let \(k\) a field, B the ring \(k[[T]]\) , \(\mathrm{K} = k((\mathrm{T}))\) its field of fractions, and A the sub- \(k\) -algebra of B generated by the elements \(\mathrm{T}^{\gamma}\) for \(\gamma \in \Gamma\) .
\end{enumerate}

\begin{enumerate}
\def\labelenumi{\alph{enumi})}
\item
  Prove that A is a local integral noetherian ring of dimension 1, whose field of fractions is K and whose integral closure is B; the ideal c = A : B is equal to BT \(^{c}\) .
\item
  For A to be a Gorenstein ring, it is necessary and sufficient that one have \(c = 2\operatorname{Card}(\mathbf{N} - \Gamma)\) .
\item
  Let \(\Omega\) the sub-A-module of K generated by the elements \(T^{-\alpha}\) for \(\alpha \notin \Gamma\) . Prove that \(K/\Omega\) is a Matlis A-module (use prop. 2 of § 8), and that \(\Omega\) is a dualizing A-module.
\end{enumerate}

\begin{enumerate}
\def\labelenumi{\arabic{enumi})}
\setcounter{enumi}{9}
\tightlist
\item
  Let A be a noetherian ring. Let I be a bounded injective A-complex, whose homology is of finite type; for every A-complex C, we denote by D(C) the complex Homgr \(_{A}\) (C,I).
\end{enumerate}

\begin{enumerate}
\def\labelenumi{\alph{enumi})}
\tightlist
\item
  Prove that the following conditions are equivalent:
\end{enumerate}

\begin{enumerate}
\def\labelenumi{(\roman{enumi})}
\item
  for every finitely generated A-module M, the canonical morphism of complexes \(\boldsymbol{\alpha}_{\mathrm{M}}:\mathrm{M}\to\mathbf{D}(\mathbf{D}(\mathrm{M}))\) (n° 5) is a homomorphism;
\item
  for every bounded A-complex C whose homology is of finite type, the morphism \(\alpha_{\mathrm{C}}: \mathrm{C} \to \mathbf{D}(\mathbf{D}(\mathrm{C}))\) (loc. cit.) is a homologism;
\item
  the canonical morphism \(\alpha_{\mathrm{A}}:\mathrm{A}\to \mathrm{Homgr}_{\mathrm{A}}(\mathrm{I},\mathrm{I})\) is a homological isomorphism.
\end{enumerate}

(Adapt the proof of Thm. 1.)

We say that the complex I is dualizing if it is bounded, injective, H(I) is finitely generated, and it satisfies the equivalent conditions above.

\begin{enumerate}
\def\labelenumi{\alph{enumi})}
\setcounter{enumi}{1}
\item
  If A admits a dualizing module \(\Omega\) any finite-length injective resolution (I, δ) of \(\Omega\) defines a dualizing complex (loc. cit.).
\item
  Let B be an A-algebra which is a finitely generated A-module, and let I be a dualizing complex of A-modules; prove that the complex of B-modules \(\mathrm{Homgr}_{\mathrm{A}}(\mathrm{B},\mathrm{I})\) is dualizing. Thus every quotient ring of a Gorenstein ring of finite dimension admits a dualizing complex.
\item
  Let I be a bounded injective A-complex. For I to be dualizing, it is necessary and sufficient that the complex of \(A_{m}\) -modules \(I_{m}\) so that it is dualizing for every maximal ideal m of A.
\item
  Let I be a dualizing A-complex, P a projective A-module of rank 1, n an integer. The complex \(1 \otimes_{\mathrm{A}} \mathrm{P}(n)\) is dualizing.
\item
  Let I, J be bounded injective A-complexes whose homology is of finite type, and \(u: I \rightarrow J\) a homomorphism; for J to be dualizing, it is necessary and sufficient that I be so.
\end{enumerate}

¶ 11) Let A be a Noetherian ring.

\begin{enumerate}
\def\labelenumi{\alph{enumi})}
\tightlist
\item
  Assume \(\Lambda\) local. Let P, Q be bounded flat A-complexes. Suppose that \(\mathrm{H}(\mathrm{P}\otimes_{\mathrm{A}}\mathrm{Q})\) is zero in degree \(\neq 0\) and free of rank 1 in degree 0. Prove that there exists an integer \(p\in \mathbf{Z}\) and homologisms \(\mathrm{A}\to \mathrm{P}(p)\) and \(\mathrm{A}\to \mathrm{Q}(-p)\) .
\end{enumerate}

(Using lemma 1 and prop. 1 of A, X, pp.~66 and 62, reduce to the case where P and Q are zero on the right; using prop. 1 of loc. cit. and th. 3 of II, § 5, no. 4, then construct complexes P' and Q' such that P = A ⊕ P' and Q = A ⊕ Q', and prove that one has H(P') = H(Q') = 0.)

\begin{enumerate}
\def\labelenumi{\alph{enumi})}
\setcounter{enumi}{1}
\item
  Let I, J be dualizing A-complexes (exerc. 10). If A is local, prove that there exists an integer n and a homologism of J onto I(n) (let P = Homgr \(_A\) (I, J), Q = Homgr \(_A\) (J, I); using the morphism \(\nu\) of no. 7 and the homologism \(\alpha_{J}\) : J → Homgr \(_A\) (Q, I) relative to the dualizing complex 1, construct a homomorphism P ⊗ \(_A\) Q → A, and apply a)).
\item
  In the general case, prove that there exists an integer n, a projective A-module L of rank 1, and a homomorphism from J onto \(I \otimes_{A} L(n)\) (set \(L = H(\text{Homgr}_{A}(I, J))\) , and apply b)).
\end{enumerate}

\begin{enumerate}
\def\labelenumi{\arabic{enumi})}
\setcounter{enumi}{11}
\tightlist
\item
  Let \(k\) be a field, and R a \(k\) -graded algebra of type \(\mathbf{N}\) such that \(\mathrm{R}_0 = k\) and that R be a Macaulay ring, of dimension \(d\) For every finitely generated graded R-module M, we denote by \(\mathrm{P}_{\mathrm{M}}(\mathrm{T})\) the Poincaré series \(\sum_{i} (\dim_k \mathrm{M}_i) \mathrm{T}^i\) and one sets \(\mathrm{Q}_{\mathrm{M}}(\mathrm{T}) = (1 - \mathrm{T})^{\dim(M)} \mathrm{P}_{\mathrm{M}}(\mathrm{T})\) ; it is an element of \(\mathbf{Z}[\mathrm{T}, \mathrm{T}^{-1}]\) (VIII, § 6, no. 3, prop. 5).
\end{enumerate}

\begin{enumerate}
\def\labelenumi{\alph{enumi})}
\item
  Let \(\Omega\) a dualizing R-module; prove that \(\Omega\) admits a graded R-module structure for which \(Q_{\Omega}(T) = (-1)^{d} Q_{\mathbb{R}}(T^{-1})\) (write R as a quotient of a polynomial algebra \(A = k[X_1, \ldots, X_n]\) , with \(\deg(X_i) = d_i\) ; if L is a free graded resolution of the A-module R, observe that the complex \(\mathrm{Homgr}_{\mathrm{A}}(\mathrm{L}, \mathrm{A})(-d)\) defines one of \(\Omega\) , and use Example 3 of VIII, § 4, no. 2).
\item
  Assume that R is a Gorenstein ring; prove that there exists an integer \(a \in \mathbf{Z}\) satisfying \(\mathrm{Q}_{\mathbb{R}}(\mathrm{T}^{-1}) = (-1)^{d}\mathrm{T}^{a}\mathrm{Q}_{\mathbb{R}}(\mathrm{T})\) .
\item
  Assume that R is an integral domain and that there exists an integer \(a \in \mathbf{Z}\) such that one has \(\mathrm{Q}_{\mathrm{R}}(\mathrm{T}^{-1}) = (-1)^{d}\mathrm{T}^{a}\mathrm{Q}_{\mathrm{R}}(\mathrm{T})\) Prove that R is a Gorenstein ring (endow \(\Omega\) of a graduation for which \(\mathrm{Q}_{\Omega} = \mathrm{Q}_{\mathrm{R}}\) , and prove that a nonzero element of \(\Omega_{0}\) forms a basis of \(\Omega\) ).
\item
  Let R be the graded k-algebra \(k[X,Y]/(X^{3},XY,Y^{2})\) ; prove that we have \(Q_{R}(T^{-1}) = T^{-2}Q_{R}(T')\) , but that R is not a Gorenstein ring.
\end{enumerate}

\exercisesection{Local cohomology, Grothendieck duality}

\begin{enumerate}
\def\labelenumi{\arabic{enumi})}
\tightlist
\item
  Let A be a ring, J an ideal of A, M an A-module. We denote by \(H_{A,J}(M)\) , or simply \(H_{J}(M)\) , the graded A-module \(\varinjlim_{n}\operatorname{Ext}_{A}(A/J^{n},M)\) . If A is local Noetherian, the graded A-module \(H_{A,m_{A}}(M)\) is identified with \(H_{A}(M)\) .
\end{enumerate}

\begin{enumerate}
\def\labelenumi{\alph{enumi})}
\item
  Establish for \(H_{J}(M)\) the properties analogous to those of \(H_{A}(M)\) with respect to homomorphisms and exact sequences of A-modules (no. 1).
\item
  Define for every \(i \geqslant 1\) isomorphisms \(\mathrm{H}_{\mathrm{J}}^{i+1}(\mathrm{M}) \to \varinjlim \operatorname{Ext}_{\mathrm{A}}^{i}(\mathrm{J}^{n}, \mathrm{M})\) .
\item
  We suppose the ring A is local Noetherian of dimension d, and the ideal J is generated by a completely secant sequence \((x_{1},\ldots,x_{r})\) . Construct an isomorphism

\[
\tau^ {i} (\mathrm{J}; \mathrm{M}): ^ {\prime} \operatorname{Tor} _ {r - i} ^ {\mathrm{A}} (\mathrm{M}, \mathrm{H} _ {\mathrm{J}} ^ {d} (\mathrm{A})) \longrightarrow \mathrm{H} _ {\mathrm{J}} ^ {i} (\mathrm{M})
\]

which generalizes the isomorphism \(\tau^{i}(\mathrm{M})\) of no. 2; in particular, we have \(\mathrm{H}_{\mathrm{J}}^{i}(\mathrm{A}) = 0\) for \(i > r\) .

\item
  For every ideal K of A contained in J, we denote by \(\rho_{\mathrm{K},\mathrm{J}}:\mathrm{H}_{\mathrm{J}}(\mathrm{M})\to \mathrm{H}_{\mathrm{K}}(\mathrm{M})\) the homomorphism deduced from the canonical maps of \(\mathrm{A / K^n}\) in \(\mathrm{A / J^n}\) . Let I be an ideal of A; construct a long exact sequence of A-modules

\[
\dots \longrightarrow \mathrm{H} _ {\mathrm{I} +, \mathrm{J}} ^ {i} (\mathrm{M}) \stackrel {\alpha} {\longrightarrow} \mathrm{H} _ {\mathrm{I}} ^ {i} (\mathrm{M}) \oplus \mathrm{H} _ {\mathrm{J}} ^ {i} (\mathrm{M}) \stackrel {\beta} {\longrightarrow} \mathrm{H} _ {\mathrm{I} \cap \mathrm{J}} ^ {i} (\mathrm{M}) \longrightarrow \mathrm{H} _ {\mathrm{I} \mid \mathrm{J}} ^ {i + 1} (\mathrm{M}) \longrightarrow \dots ,
\]

\[
\text {   où   } \alpha (x) = (\rho_ {\mathrm{I,I+J}} (x), \rho_ {\mathrm{J,I+J}} (x)) \text {   et   } \beta (x, y) = \rho_ {\mathrm{I} \cap \mathrm{J}, \mathrm{I}} (x) - \rho_ {\mathrm{I} \cap \mathrm{J}, \mathrm{J}} (y).
\]

\end{enumerate}

\begin{enumerate}
\def\labelenumi{\arabic{enumi})}
\setcounter{enumi}{1}
\tightlist
\item
  Let \(\Lambda\) a local Macaulay ring of dimension \(d\) . Prove that for A to be a Gorenstein ring, it is necessary and sufficient that the \(\Lambda\) -module \(\mathrm{H}_{\mathrm{A}}^{d}(\Lambda)\) be injective.
\end{enumerate}

¶ 3) Let \(\rho: A \to B\) be a local homomorphism of local Noetherian rings, such that \(\rho(\mathfrak{m}_{A})B\) is a definition ideal of B.

\begin{enumerate}
\def\labelenumi{\alph{enumi})}
\item
  Construct for every B-module N a natural A-linear isomorphism of \(H_{A}(N)\) on \(H_{B}(N)\) . (Using an injective resolution of N, show that it suffices to prove that we have \(H_{A}^{i}(J)=0\) for all i\textgreater0 and every indecomposable injective B-module J. Let q be the unique element of \(\operatorname{Ass}_{\mathrm{B}}(\mathrm{J})\) , and by \(p=\rho^{-1}(q)\) If \(p\neq m_{A}\) , observe that there exists \(a\in m_{A}\) such that the homothety \(a_{J}\) be bijective. If \(p=m_{A}\) , and if I is a minimal injective resolution of the A-module J, observe that we have \(\operatorname{Ass}_{\mathrm{A}}(\mathrm{I}^{n})=\{\mathfrak{m}_{\mathrm{A}}\}\) for every n.)
\item
  We suppose moreover that B is a flat A-module. Construct for every A-module M a natural B-linear isomorphism of \(\mathrm{B}\otimes_{\mathrm{A}}\mathrm{H}_{\mathrm{A}}(\mathrm{M})\) on \(\mathrm{H}_{\mathrm{B}}(\mathrm{B}\otimes_{\mathrm{A}}\mathrm{M})\) .
\end{enumerate}

\begin{enumerate}
\def\labelenumi{\arabic{enumi})}
\setcounter{enumi}{3}
\item
  Let A be a local ring and M a nonzero finitely generated A-module, of dimension e. Prove that \(\mathrm{H}_{\mathrm{A}}^{e}(\mathrm{M})\) is not zero (reduce to the case where the ring A is complete, then with the aid of exerc. 3 to the case where A is regular).
\item
  Let I be a finite subset of N.
\end{enumerate}

\begin{enumerate}
\def\labelenumi{\alph{enumi})}
\item
  Construct a regular local ring B and a B-module N such that the set of centers p such that \(H_{\mathrm{B}}^{p}(\mathrm{N}) \neq 0\) is equal to I (consider a direct sum).
\item
  Construct a local ring A such that the set of integers p such that \(\mathrm{H}_{\mathrm{A}}^{p}(\mathrm{A}) \neq 0\) is equal to I (take \(A = B \oplus N\) , where N is an ideal of square zero, and use exerc. 3).
\end{enumerate}

\begin{enumerate}
\def\labelenumi{\arabic{enumi})}
\setcounter{enumi}{5}
\tightlist
\item
  Let A be a local Noetherian ring admitting a dualizing module, M a finitely generated A-module, n an integer.
\end{enumerate}

\begin{enumerate}
\def\labelenumi{\alph{enumi})}
\item
  For the A-module \(H_{\mathrm{A}}^{n}(\mathrm{M})\) to be of finite length, it is necessary and sufficient that for every prime ideal p of A distinct from \(m_{A}\) , one has \(\mathrm{H}_{\mathrm{A}_{\mathfrak{p}}}^{n-\dim(\mathrm{A}/\mathfrak{p})}(\mathrm{M}_{\mathfrak{p}})=0\) (use Grothendieck duality).
\item
  For \(\mathrm{H}_{\mathrm{A}}^{i}(\mathrm{M})\) be of finite length for every \(i \leqslant n\) , it is necessary and sufficient that one has \(\operatorname{prof}(\mathrm{M}_{\mathfrak{p}}) + \dim(\mathrm{A}/\mathfrak{p}) > n\) for every prime ideal p distinct from \(m_{A}\)
\end{enumerate}

\begin{enumerate}
\def\labelenumi{\arabic{enumi})}
\setcounter{enumi}{6}
\tightlist
\item
  Let A be a ring, M a \(\Lambda\) -module, \(\mathbf{x} = (x_i)_{i \in I}\) a finite family of elements of A. For every integer \(n \geqslant 1\) , we denote \(\mathbf{x}^n\) the family \((x_i^n)_{i \in I}\) . Let \(\Delta(\mathbf{x})\) the endomorphism of \(\mathrm{A}^{\mathrm{I}}\) defined by \(\Delta(\mathbf{x})(e_i) = x_i e_i\) .
\end{enumerate}

\begin{enumerate}
\def\labelenumi{\alph{enumi})}
\tightlist
\item
  For every pair of integers \((n, m)\) such that \(n \geqslant m\) show that the maps
\end{enumerate}

\[
\operatorname{Homgr} \left(\Lambda \left(\Delta \left(\mathbf {x} ^ {n - m}\right)\right), 1 _ {\mathrm{M}}\right): \mathbf {K} ^ {\bullet} \left(\mathbf {x} ^ {m}, \mathrm{M}\right) \longrightarrow \mathbf {K} ^ {\bullet} \left(\mathbf {x} ^ {n}, \mathrm{M}\right)
\]

define an inductive system of complexes; we denote by \(\mathbf{K}^{\bullet}(\mathbf{x}^{\infty},\mathrm{M})\) the limit of this system, and \(\mathrm{H}^{\bullet}(\mathbf{x}^{\infty},\mathrm{M})\) its cohomology.

\begin{enumerate}
\def\labelenumi{\alph{enumi})}
\setcounter{enumi}{1}
\tightlist
\item
  We choose a total order on I. For every subset J of I, we set \(x_{J} = \prod_{j \in J} x_{j}\) .
\end{enumerate}

For \(p \in \mathbf{Z}\) , we set \(\mathcal{C}^p = \bigoplus_{|J| = p} A_{x_J}\) ; for \(a \in A_{x_J}\) , with \(\text{Card}(J) = p\) , we set \(d^p(a) = \sum_{i \notin J} (-1)^{\varepsilon(J, i)} \frac{a}{1}\) in \(A_{x_{J \cup \{i\}}}\) , where \(\varepsilon(J, i)\) is the number of elements of \(J\) strictly less than \(i\) Prove that \((\mathcal{C}, d)\) is a complex, and define a canonical isomorphism of \(\mathbf{K}^\bullet(\mathbf{x}^\infty, M)\) on the complex \(\mathcal{C} \otimes_A M\) .

8) We keep the notation from the previous exercise, assuming moreover that the ring A is Noetherian. We denote by x the ideal generated by x.

\begin{enumerate}
\def\labelenumi{\alph{enumi})}
\item
  Let m be an integer. Prove that there exists an integer \(n_{0} \geqslant m\) such that for \(n \geqslant n_{0}\) such that the morphism \(\boldsymbol{\Lambda}(\Delta(\mathbf{x})^{n-m}) : \mathbf{K}_{\bullet}(\mathbf{x}^{n}, \mathrm{A}) \longrightarrow \mathbf{K}_{\bullet}(\mathbf{x}^{m}, \mathrm{A})\) induces 0 on the homology in degree \(\geqslant 1\) (observe that one has \(\boldsymbol{\Lambda}^{p}(\Delta(\mathbf{x}^{n-m}))(\mathbf{K}_{p}(\mathbf{x}^{n}, \mathrm{A})) \subset x^{n-m}\mathbf{K}_{p}(\mathbf{x}^{m}, \mathrm{A})\) ; deduce from the Artin–Rees lemma that there exists an integer \(n_{0}\) such that \(Z_{p}(\mathbf{x}^{m}, \mathrm{A}) \cap x^{n_{0}}\mathbf{K}_{\bullet}(\mathbf{x}^{m}, \mathrm{A})\) is contained in \(x^{rm}Z_{p}(\mathbf{x}^{m}, \mathrm{A})\) , and hence in \(B_{p}(\mathbf{x}^{m}, \mathrm{A})\) ).
\item
  Define a canonical isomorphism from \(\varinjlim_{n}\mathrm{Ext}_{\mathrm{A}}^{p}(\mathrm{A}/x^{n},\mathrm{M})\) on \(\mathrm{H}^p (\mathbf{x}^\infty ,\mathrm{M})\) (reduce to proving that \(\mathrm{H}^p (\mathbf{x}^\infty ,\mathrm{M})\) is zero for \(p\geqslant 1\) and M injective; in this case observe that \(\mathrm{H}^p (\mathbf{x}^\infty ,\mathrm{M})\) is identified with the inductive limit of the \(\mathrm{Hom}_{\mathrm{A}}(\mathrm{H}_p(\mathbf{x}^n,\mathrm{M}))\) and use \(a)\) ). In particular, if A is local and if \(x\) is an ideal of definition of A, the graded A-module \(\mathrm{H}_{\mathrm{A}}(\mathrm{M})\) is canonically isomorphic to \(\mathrm{H}^{\bullet}(\mathbf{x}^{\infty},\mathrm{M})\) .
\item
  We assume the ring A is local Noetherian, of dimension \(d\) ; let \(\mathbf{x} = (x_{1},\ldots ,x_{d})\) a maximal secant sequence of elements of \(\mathfrak{m}_{\Lambda}\) Show that the A-module \(\mathrm{H}_{\mathrm{A}}^{d}(\mathrm{M})\) is identified with the limit of the inductive system of A-modules \((\mathrm{M}_p,u_{qp})\) , where \(\mathrm{M}_p\) is the quotient module \(\mathrm{M} / (x_1^p\mathrm{M} + \dots +x_d^p\mathrm{M})\) and \(u_{qp}:\mathrm{M}_p\to \mathrm{M}_q\) ( \(q\geqslant p\) ) the homomorphism induced by the homothety of ratio \((x_{1}\dots x_{d})^{q - p}\) .
\end{enumerate}

\begin{enumerate}
\def\labelenumi{\arabic{enumi})}
\setcounter{enumi}{8}
\tightlist
\item
  Let A be a Noetherian local ring, d its dimension, \(\mathbf{x} = (x_{1}, \ldots, x_{d})\) a maximal secant sequence of elements of \(m_{A}\) .
\end{enumerate}

\begin{enumerate}
\def\labelenumi{\alph{enumi})}
\item
  Prove that there exists an integer \(m_0\) such that for \(m \geqslant m_0\) and \(n \geqslant 0\) , one has \((x_1^m \ldots x_d^m)^n \not\in (x_1^{m(n+1)}, \ldots, x_d^{m(n+1)})\) (use Exercises 4 and 7, b)).
\item
  Assume that \(\Lambda\) contains a field of characteristic \(p\) . Prove that we have \((x_{1} \ldots x_{d})^{n} \notin (x_{1}^{n+1}, \ldots, x_{d}^{n+1})\) for every \(n \geqslant 0\) , that is, A satisfies condition (CM) (Exercise 3 of § 2; consider the endomorphism \(x \mapsto x^p\) Therefore, every Noetherian local ring containing a field satisfies condition (CM). \(^{1}\) .
\item
  We further assume that A is regular; let \(\rho: A \to B\) an injective ring homomorphism making B into a finitely generated A-module. Prove that the A-module \(\rho(A)\) is a direct factor in B (cf.~§ 2, exerc. 3); reduce to the case where B is integral, hence local, and apply exerc. 8 of § 4).
\end{enumerate}

\begin{enumerate}
\def\labelenumi{\arabic{enumi})}
\setcounter{enumi}{9}
\tightlist
\item
  Let \(k\) a field, A a \(k\) -regular local algebra of dimension \(d\) , such that the extension \(k \to \kappa_{\mathrm{A}}\) is trivial. The A-modules \(\mathrm{H}_{\mathrm{A}}^{d}(\mathrm{A})\) and \(\mathrm{A}' = \mathrm{Hom}_k^{cont}(\mathrm{A}, k)\) are Matlis modules, hence isomorphic; we propose to make explicit an isomorphism between these two modules.
\end{enumerate}

\begin{enumerate}
\def\labelenumi{\alph{enumi})}
\item
  Let \(\mathbf{x} = (x_{1}, \ldots, x_{d})\) a coordinate system of A. For every integer \(p \geqslant 0\) , we denote \(A_{p}\) the k-algebra \(\mathrm{A}/(x_{1}^{p}, \ldots, x_{d}^{p})\) The classes of monomials \(x_{1}^{a_{1}} \ldots x_{d}^{a_{d}}\) with \(0 \leqslant a_{i} < p\) form a basis of \(A_{p}\) on k; we denote by \(\rho_{p}\) the linear form on \(A_{p}\) which takes the value 1 on the class of \(x_{1}^{p-1} \ldots x_{d}^{p-1}\) and 0 on the other elements of the basis. Prove that the k-linear map \(\tilde{\rho}_{p}: A_{p} \to \operatorname{Hom}_{k}(A_{p}, k)\) derived from the bilinear form \((a, b) \mapsto \rho_{p}(ab)\) is an isomorphism (one may reason directly, or use Exercise 10 of § 5).
\item
  For \(q \geqslant p\) let \(u_{qp}: A_p \to A_q\) the A-linear homomorphism induced by the homothety of ratio \((x_1 \ldots x_d)^{q-p}\) , and let \(\pi_{pq}: A_q \to A_p\) the passage-to-the-quotient map; prove that one has \(\tilde{\rho}_q \circ u_{qp} = {}^t \pi_{pq} \circ \tilde{\rho}_p\) . Deduce by passage to the inductive limit, taking into account exerc. 8, c), an A-linear isomorphism \(\tilde{\rho}: H_A^d(A) \to A'\) .
\item
  Assume \(d=1\) , so that A is a discrete valuation ring, and \(x_{1}=x\) a uniformizer of A. Show that the choice of the uniformizer \(x\) allows one to identify \(\mathrm{H}_{\mathrm{A}}^{1}(\mathrm{A})\) to K/A; the isomorphism \(\rho:\mathrm{K}/\mathrm{A}\to\mathrm{A}'\) is then defined by \(\rho(\varphi)(a)=\mathrm{Res}(a\varphi)\) , where \(\mathrm{Res}:\mathrm{K}/\mathrm{A}\to k\) is the projection onto the factor \(k x^{-1}\) in the canonical decomposition \(\mathrm{K}/\Lambda=\bigoplus_{n\geqslant 1}k x^{-n}\) .
\end{enumerate}

\begin{enumerate}
\def\labelenumi{\arabic{enumi})}
\setcounter{enumi}{10}
\item
  Let A be a Noetherian local ring, I a Matlis module over A, and M a finitely generated A-module. We set \(\mu^{p} = \dim_{\kappa_{A}} \operatorname{Ext}_{A}^{p}(\kappa_{A}, M)\) for every \(p \geqslant 0\) (cf. § 8, exerc. 4). Prove that the graded A-module H \(_{A}\) (M) is identified with the homology of a complex \(\ldots \to 0 \to I^{\mu^{0}} \to I^{\mu^{1}} \to \ldots\) (loc. cit.). In particular, if M is nonzero, we have \(\mu^{p} \neq 0\) for \(p = \text{prof}(M)\) and \(p = \dim(M)\) (exercise 4).
\item
  Let A be a Noetherian local ring, M a finitely generated A-module of dimension d; suppose that M is a Buchsbaum module (§ 2, exerc. 7).
\end{enumerate}

\begin{enumerate}
\def\labelenumi{\alph{enumi})}
\item
  Prove that \(\mathrm{H}_{\mathrm{A}}^{i}(\mathrm{M})\) is annihilated by \(m_{A}\) for \(i \neq d\) (First handle the case d = 1, then reason by induction on d).
\item
  Let \(x\) a secant element for \(\mathfrak{m}\) ; define for \(0 \leqslant i \leqslant d - 2\) the exact sequences \(0 \to \mathrm{H}_{\mathrm{A}}^{i}(\mathrm{M}) \to \mathrm{H}_{\mathrm{A}}^{i}(\mathrm{M}/x\mathrm{M}) \to \mathrm{H}_{\mathrm{A}}^{i+1}(\mathrm{M}) \to 0\) .
\item
  Prove the formula \(i(\mathrm{M}) = \sum_{i=0}^{d-1} \left( \begin{array}{cc} d & 1 \\ i & \end{array} \right) \lg(\mathrm{H}_{\mathrm{A}}^i(\mathrm{M}))\) (argue by induction on the integer \(d \geqslant 1\) , using \(b\) ) and exerc. 8, \(d\) ) of § 2).
\end{enumerate}

\specialchapter{Index of notations}

\(\mathrm{prof}_{\mathrm{A}}(\mathrm{J};\mathrm{M})\) , \(\mathrm{prof}(\mathrm{J};\mathrm{M})\) , \(\mathrm{prof}_{\mathrm{A}}(\mathrm{M})\) , \(\mathrm{prof}(\mathrm{M})\) : p.~2. \(\mathbf{K}^{\bullet}(\mathbf{x},\mathbf{M})\) : p.~5. \(\mathrm{H}^{\bullet}(\mathbf{x},\mathbf{M})\) p.~5. \(\mathrm{prof}_{\mathrm{F}}(\mathrm{M})\) p.~10. \(\mathrm{gradc}(\mathrm{N})\) p.~11. \(\mathrm{dp}_{\mathrm{A}}(\mathrm{M})\) , \(\mathrm{di}_{\mathrm{A}}(\mathrm{M})\) p.~37. \(\mathrm{dh(A)}\) p.~38. \(\delta (\Lambda)\) p.~60. \(\alpha_{\mathrm{J}},\beta_{\mathrm{J}}\) p.~63. \(\Omega_k(\mathrm{A}),d_{\mathrm{A / k}}\) p.~78. \(\Omega (\mathrm{A})\) p.~93. \(\mathcal{I}(\mathrm{A})\) p.~105.\\
I(p), \(e_{\mathfrak{p}}\) p.~105. \(\mathrm{D_A(M),D_A(f),}\alpha_{\mathrm{M}}\) p.~112. \(\theta (\mathrm{M},\mathrm{P}),\rho (\mathrm{M},\mathrm{P})\) p.~121. \(\mathbf{D}(\mathbf{C})\) , p.~134. \(\mathbf{D}_{\mathrm{A}}(f),\mathbf{D}(\mathbf{M})\) p.~135. \(\alpha_{\mathrm{M}}\) p.~136. \(\mathrm{H}_{\mathrm{A}}(\mathrm{M}),\mathrm{H}_{\mathrm{A}}(f)\) p.~142. \(\mathcal{D}\) p.~142. \(\mathcal{D}_{cs}\) p.~145. \(\tau^i (\mathbf{M})\) p.~146. \(\gamma^i (\mathbf{M}),\delta^i (\mathbf{M})\) p.~148. \(S_{k}\) p.~151, exerc. 8. \(t_\mathrm{A}(\mathbf{M}),t(\mathbf{M})\) p.~153, exerc. 16. \(\chi (\mathbf{x},\mathbf{N})\) p.~155, exerc. 6. \(\chi (\mathbf{M})\) p.~157, exerc. 2. \(rg(u),\mathfrak{d}(u)\) p.~159, exerc. 8. \(h_i(\mathbf{A})\) p.~166, exerc. 4. \(\mu^i (\mathfrak{p},\mathbf{M})\) p.~174, exerc. 4. \(\mathrm{H}_{\mathrm{J}}(\mathrm{M})\) p.~179, exerc. 1. \(\mathbf{K}^{\bullet}(\mathbf{x}^{\infty},\mathbf{M}),\mathrm{H}^{\bullet}(\mathbf{x}^{\infty},\mathbf{M})\) p.~180, exerc. 7.

\specialchapter{Terminological index}

Absolutely regular, normal (algebra): p.~75. Absolutely regular, absolutely normal algebra: p.~75. Essentially of finite type algebra: p.~71. Formally étale, formally unramified algebra: p.~170, exerc. 1 Formally smooth algebra: p.~85. Smooth algebra: p.~104. Complete intersection ring: p.~65. Gorenstein ring: p.~48. Grothendieck ring: p.~169, exerc. 5. Locally integral ring: p.~16. Macaulay ring (or Macaulay ring): p.~30. Normal ring: p.~16. Presentable ring: p.~58. Regular ring: p.~55. Auslander-Buchsbaum (theorem of): p.~45.

Bass (theorem of): p.~49. Buchsbaum (module of): p.~155, exerc. 7.

Cohen (theorem of): p.~88. Local cohomology (module of): p.~142. Completely secant (ideal): p.~62. Dualizing complex: p.~177, exerc. 10 Canonical components of a ring: p.~15.

Different (ideal): p.~167, exerc. 7. Homological dimension of a ring: p.~38. Projective or injective dimension of a module: p.~37. Dualizing (complex): p.~177, exerc. 10. Dualizing (module): p.~125. Grothendieck duality: p.~149.

Essentially of finite type (algebra): p.~71. Essentially generating (family): p.~71.

Essentially generating family: p.~71. Formally étale, formally unramified (algebra): p.~170, exerc. 1. Formally smooth (algebra): p.~85. Strongly secant (subset): p.~28.

Gorenstein (ring, of): p.~48. Grade of a module: p.~11. Grothendieck (ring, of): p.~169, exerc. 5. Grothendieck (duality of): p.~149.

Hartshorne (theorem of): p.~17. Hilbert-Burch (theorem of): p.~160, exerc. 11. Homological (dimension): p.~38. Completely secant ideal: p.~62. Injective (dimension): p.~37. Complete intersection (ring of): p.~65. Smooth (formally smooth algebra): p.~85. Smooth (algebra): p.~104. Macaulay (ring of): p.~30. Macaulay-Cohen (criterion of): p.~31. Macaulay-Cohen (theorem of): p.~30. Macaulayan (module): p.~30. Matlis (module of): p.~110. Buchsbaum module: p.~155, exerc. 7. Local cohomology module: p.~142. Matlis module: p.~110. Dualizing module: p.~125. Macaulayan (or Macaulay) module: p.~23. Perfect module: p.~158, exerc. 4. Pure module: p.~27. Normal (ring): p.~16. Perfect (module): p.~158, exerc. 4. Strongly secant part: p.~28. Presentable (ring): p.~58. Depth of a module along a closed subset: p.~10. Depth of a module, of a local ring: p.~2. Projective (dimension): p.~37. Property \(S_k\) : p.~151, exerc. 8. n-th symbolic power of a prime ideal: p.~109. Pure (module): p.~27. Regular (ring): p.~55. Lifting of a homomorphism of algebras: p.~83. Secant (completely secant ideal): p.~62. Serre (theorems of): p.~20 and p.~54. Socle of a module: p.~111. M-regular sequence: p.~8. Auslander-Buchsbaum theorem: p.~45. Bass theorem: p.~49. Cohen theorem: p.~88. Hartshorne theorem: p.~17. Hilbert-Burch theorem: p.~160, exerc. 11. Macaulay-Cohen theorem: p.~30. Serre theorems: p.~20 and p.~54. Zariski theorems: p.~101 and p.~102. Zariski (theorems of): p.~101 and p.~102.
